\chapter{推免资格、名额分配与成绩排名}
\label{ch02}
\ist{资格与排名}
\begin{keybox}
保研里最容易吃亏的，就是这一章。推免资格不是申请了就有：学院按综合排名分配名额，学校审定后报教育部备案，而排名的算法，教育部层面并不统一规定——加权平均分还是学分绩点、含不含选修课、重修算首考还是按60分、计入前五学期还是前六学期、挂几门就出局、综测公式怎么加权，各校各院差异极大，跨校套用几乎一定出错。本章把这些问题拆成十三个角度：排名口径、排名时点、转专业、名额分配、挂科、综测、公示异议、辅修双学位、课程范围、成绩折算、排名证明、交换学分、特殊班级。每一条都遵循同一写法：先给结论，再给依据，并尽量并列两到三所学校的真实条文作对照。核心方法不是背规则，而是学会去本院《推免实施细则》里找三样东西——\textbf{课程计算清单、综合成绩公式、资格与名额的基数口径}。读完你应当能做到两件事：知道该问学院哪几个问题；拿到公示名单时，能逐项把自己的分数算回去。
\end{keybox}

\section{挂科}
\begin{qa}{M0069}{挂了科但补考过了，还算不算挂科记录}
\textbf{补考通过不等于没有记录。} 教育部令第41号《普通高等学校学生管理规定》第十八条明确要求：学校“对通过补考、重修获得的成绩，应当予以标注”[\nref{S0040}/\nref{S0366}]。所以你的成绩单上一定会显示这门课是补考/重修取得的。

但这门“挂过的课”算不算影响推免，各校口径完全不同，常见三种：

\begin{itemize}
  \item \textbf{只认最终状态}：要求“无课程不及格记录”，补考通过即视为满足（如广东外语外贸大学2026年办法要求“无课程不及格记录”\nref{S0134}）。
  \item \textbf{按合格线计分}：中南大学材料科学与工程学院2027年细则规定“补考和重修合格的课程成绩按60分计算”，且“有不及格课程者不得推荐”，但全校性选修课已修满6学分且及格的不影响资格\nref{S0273}。
  \item \textbf{只看首考成绩}：成都大学商学院、北京航空航天大学计算机学院均规定课程成绩“采用首次获得的成绩”\nref{S0079}\nref{S0579}。
\end{itemize}

结论：记录一定留下，能不能过资格审查要看本校写的是“无不及格记录”还是“首考成绩”。
\nref{S0040}
\nref{S0366}
\nref{S0134}
\nref{S0273}
\nref{S0079}
\nref{S0579}
\end{qa}

\iin{中华人民共和国教育部}
\ith{挂科}
\ist{资格与排名}
\begin{qa}{M0070}{挂科超过几门会取消推免资格}
\textbf{教育部层面没有规定“挂几门取消资格”的红线}，各校自设门槛，差异从“0门”到“2门”都有：

\par\vspace{2pt}\nobreak\begin{center}\footnotesize
\begin{tabularx}{\textwidth}{>{\raggedright\arraybackslash}X>{\raggedright\arraybackslash}X}
\toprule
\textbf{学校（依据年份）} & \textbf{挂科口径} \\ \midrule
广东外语外贸大学（2026年办法） & “无课程不及格记录”，等于0门\nref{S0134} \\
大连理工大学商学院 & “一考不及格的必修课程不能超过2门”，所有参评课程以一考成绩计算\nref{S0549} \\
云南师范大学（2027年通知） & 专业必修、专业限选课\textbf{无}不合格记录；其他各类型课程不合格记录不超过2门\nref{S0568} \\
中南大学材料学院（2027年细则） & “有不及格课程者不得推荐”；全校性选修课已修满6学分且及格的不影响资格\nref{S0273} \\
温州医科大学（2027年通知） & 必修课及限选课无不及格，且平均学分绩点≥4.0\nref{S0563} \\
\bottomrule
\end{tabularx}
\end{center}\vspace{2pt}\par

注意两件事：一是很多学校卡的是“\textbf{一考}成绩”，补考通过也不豁免（大连理工\nref{S0549}）；二是门槛之外还有排序，即使门数没超，低分课程也会拉低加权分。以本校当年文件为准。
\nref{S0134}
\nref{S0549}
\nref{S0568}
\nref{S0273}
\nref{S0563}
\end{qa}

\iin{广东外语外贸大学}
\ith{挂科}
\ist{资格与排名}
\begin{qa}{M0072}{竞赛加分只认国奖吗？省赛一等奖算不算}
\textbf{没有全国统一答案}，各校对“什么级别的竞赛能加分”写得差异极大，必须查本校当年的加分目录：

\begin{itemize}
  \item \textbf{只认国家级及以上}：湖南大学2027年推免创新创业赛事加分公示中，只公示“国家级及以上奖励”，如挑战杯国家级特等奖排名第1加6分、国家级三等奖排名第1加2分、国家级三等奖排名第4加0.67分\nref{S0541}。
  \item \textbf{国家级＋省级都算}：中南大学材料科学与工程学院2027年细则规定，国家级重点赛事一等奖排名第一加0.5分；“教育部等中央部委、省教育厅、省科技厅组织的重点赛事以及学校重点赛事对应的省级比赛”一等奖排名第一同样加0.5分\nref{S0273}。
  \item \textbf{只认全国赛且排除省级选拔赛}：安徽师范大学2027年方案写明“全国赛不含省级赛或国家级赛事省级选拔赛和区域性比赛”，且仅限其竞赛库中B2及以上项目\nref{S0646}。
  \item \textbf{只认全国赛一等奖}：内蒙古大学规定作为主力参加国内权威科研竞赛（全国赛）获一等奖，学分绩点加分最高不超过0.15\nref{S0544}。
\end{itemize}

所以省赛一等奖在有的学校能加分、分值甚至与国奖相同，在有的学校一分不加。关键是本院细则里的\textbf{竞赛认定目录}，且多数学校要求同一类别多项只取最高、不累计。
\nref{S0541}
\nref{S0273}
\nref{S0646}
\nref{S0544}
\nref{S0533}
\end{qa}

\iin{湖南大学}
\ith{竞赛加分}
\ist{资格与排名}
\begin{qa}{M0094}{少数民族预科班转正后算不算应届推免生}
\textbf{现行公开文件中未见针对“少数民族预科班转正后是否算应届推免生”的全国统一规定}，须以本校当年文件为准。

可以确认的相关规则有两条：

\begin{enumerate}
  \item \textbf{推免对象的一般定义}。教育部《全国普通高等学校推荐优秀应届本科毕业生免试攻读硕士学位研究生工作管理办法（试行）》（教学〔2006〕14号）第十四条要求推荐对象为“纳入国家普通本科招生计划录取的应届毕业生”\nref{S0454}。预科班转入本科后，只要学籍已转为普通本科、且推免当年正处于毕业学年，形式上符合这一表述——但“毕业学年”的认定由学校学籍管理部门作出。
  \item \textbf{少数民族学生的另一条独立通道}。“少数民族高层次骨干人才计划”是\textbf{研究生招生}阶段的专项，考生须在研招网骨干计划资格申请平台提交申请并经审核，毕业后按协议回定向地区就业\nref{S0573}；它与本科推免资格是两回事，不能互替。
\end{enumerate}

实操上你只需确认两点：本校实施办法的“推荐对象”条款有没有排除预科转入学生；以及教务系统里你的学籍状态是否显示为应届毕业学年。
\nref{S0454}
\nref{S0573}
\nref{S0079}
\end{qa}

\iin{教育部学生服务与素质发展中心（学信网/研招网）}
\ith{预科}
\ist{资格与排名}
\begin{qa}{M0097}{辅导员专项思政骨干专项是怎么报的，算不算推免}
要分清楚：\textbf{教育部层面没有面向本科应届生的“辅导员专项”“思政骨干专项”推免通道。}

\begin{itemize}
  \item \textbf{思政骨干专项}指的是“高校思想政治工作骨干在职攻读博士学位专项计划”，面向的是\textbf{已在职}的高校思想政治工作骨干，不是在读本科生，与本科推免不是一回事。
  \item 教育部关于研究生的“三助一辅”文件（《关于做好研究生担任助研、助教、助管和学生辅导员工作的意见》，2014年）说的是研究生在读期间担任学生辅导员的\textbf{岗位}安排，属于培养环节，也不是推免渠道\nref{S0053}。
\end{itemize}

但\textbf{部分高校确实在校内推免专项里设“辅导员计划”}。例如东南大学在夏令营通知中要求：本校应届本科生拟报名“工程硕博士专项、集成电路专项班、国优计划、支教团计划、辅导员计划等推免专项”的，须报名参加夏令营\nref{S0126}。也就是说，这类专项由学校自主设置，通常做法是：先按普通程序取得推免资格，再报名参加该专项的选拔；名额来源、是否单列、是否限定专业，全看本校文件。

结论：想报哪个专项，直接查本校当年“推免专项报名通知”，不要按名字去套全国政策。
\nref{S0053}
\nref{S0126}
\nref{S0512}
\nref{S0567}
\end{qa}

\iin{中华人民共和国教育部}
\ith{专项计划}
\ist{资格与排名}
\begin{qa}{M0098}{创新学分不够会不会影响推免资格}
\textbf{看本校有没有把创新学分写成硬条件——写了的，就是硬门槛。}

“创新创业学分”本身是学校自设的制度。教育部令第41号第十七条规定：学生参加创新创业、社会实践等活动以及发表论文、获得专利授权等与专业学习、学业要求相关的经历成果，“可以折算为学分，计入学业成绩。具体办法由学校规定”，学校“可以建立创新创业档案、设置创新创业学分”\nref{S0040}。所以有没有、要多少，各校自定。

已经明文写进推免条件的：
\begin{itemize}
  \item 燕山大学2027年通知把“满足各专业创新创业教育学分要求”列为遴选条件之一\nref{S0533}；
  \item 成都大学商学院要求“修读并获得本专业人才培养方案规定的前三学年（五年制为前四学年）全部应修学分”，因教学计划调整未开设的经学院认定可视为满足\nref{S0079}；
  \item 南开大学要求“在推荐遴选时已经按照正常教学计划修满必修课程”\nref{S0561}；
  \item 内蒙古大学要求“按照专业培养方案前三学年应修的必修课均已修完，所修培养方案内课程均已合格”\nref{S0544}。
\end{itemize}

结论：如果创新学分属于你校培养方案的应修学分，没修满就可能卡在“学分条件”上；如果不属于，则一般不单独设置门槛。
\nref{S0533}
\nref{S0079}
\nref{S0561}
\nref{S0544}
\nref{S0040}
\end{qa}

\iin{zsjyc（高校官网）}
\ith{创新学分}
\ist{资格与排名}
\begin{qa}{M0099}{体测不及格会不会取消保研资格}
\textbf{没有全国统一的“体测不及格即取消推免”规定，但不少学校把体测/体育课写进了基本条件。}

\begin{itemize}
  \item 北京科技大学2027年推免通知的基本条件里明确“身心健康，体育课成绩合格”\nref{S0233}；
  \item 中山大学2026年强基计划本研转段方案要求“国家规定的学生体质健康测试达标（免测除外）”，这是一条明确的达标型硬条件\nref{S0248}；
  \item 安徽师范大学2027年方案把“体质健康测试成绩”列入学院排名公布内容、公示内容，并要求9月1日前完成“体质健康达标等情况复核”，由体委负责核定成绩\nref{S0646}。
\end{itemize}

制度依据方面，教育部令第41号第十四条规定“学生体育成绩评定要突出过程管理，可以根据考勤、课内教学、课外锻炼活动和体质健康等情况综合评定”\nref{S0040}——即如何评定由学校规定。

所以你的问题答案是：\textbf{没写就一般不影响；写了“体育课合格”或“体测达标”，就必须先达标再谈排名。} 请翻本校推免条件的“基本条件”一节逐条核对。
\nref{S0233}
\nref{S0248}
\nref{S0646}
\nref{S0040}
\end{qa}

\iin{北京科技大学}
\ith{体测}
\ist{资格与排名}
\begin{qa}{M0101}{学术不端被记录后还能不能推免}
\textbf{基本是一票否决，而且“记录”本身就是门槛。} 这条几乎全国一致。

\textbf{全国层面}：教育部办公厅关于规范和加强推免工作的通知要求，对在推免过程中弄虚作假、有论文抄袭或虚报获奖成果等学术不端行为的学生，“一经查实，取消推免资格，已入学的应当取消学籍”，并由省级教育招生考试机构按规定记入《国家教育考试考生诚信档案》\nref{S0527}\nref{S0318}。各校转述后写入本校办法，如南开大学2021年实施办法第二十一条、北京大学2027年推免办法第7条均为此表述\nref{S0561}\nref{S0553}。

\textbf{校内条件层面}：各校的“基本条件”普遍写“无任何考试作弊和学术不端行为记录”或“无任何学术不端行为记录”。例如西安电子科技大学2027年接收条件第5条“无任何学术不端行为记录，无任何不良诚信或违法违纪行为记录”\nref{S0157}；武汉轻工大学2027年办法要求“无任何考试作弊和剽窃他人学术成果等学术不端行为”\nref{S0409}。

\textbf{认定标准}：《高等学校预防与处理学术不端行为办法》（教育部令第40号）第二十七条列举了剽窃抄袭、篡改成果、伪造数据、不当署名等构成学术不端的行为\nref{S0042}。

结论：只要被正式认定为学术不端并形成记录，推免通道基本关闭；即便未形成处分，若推免材料被判“有名无实”，也会被取消资格。
\nref{S0527}
\nref{S0318}
\nref{S0561}
\nref{S0553}
\nref{S0157}
\nref{S0409}
\nref{S0042}
\end{qa}

\iin{教育部学生服务与素质发展中心（学信网/研招网）}
\ith{学术不端}
\ist{资格与排名}
\begin{qa}{M0117}{拿到推免资格但不去读会不会被记入诚信档案}
要分三种情形，结论完全不同。\textbf{单纯放弃一般不记诚信档案，但代价在别处。}

\begin{enumerate}
  \item \textbf{校内公示前主动放弃}：多数学校允许，但要求书面留痕。山东中医药大学中医学院2027年细则规定“放弃推免名额的学生，须签订书面的放弃推免生名额承诺书”\nref{S0639}；安徽师范大学2027年方案要求“学生本人及时提交书面申明、家长知情签字，明确表述自愿放弃，结果自负”\nref{S0646}。
  \item \textbf{已被接收单位录取后放弃}：几乎不被允许。教育部教学〔2006〕14号第十三条：截至教育部规定时间后，被确定的推免生无招生单位接收或本人放弃，“该生推免生资格作废，名额不得转让”\nref{S0454}。湘潭大学2026年通知更直接：“凡是被本校或外校接收的被推荐者，学校将不受理其放弃免试攻读研究生学位的申请”\nref{S0567}。河北大学2027年办法要求推荐高校对未达学业要求或发生严重违纪的推免生取消资格并通报\nref{S0604}。
  \item \textbf{记入《国家教育考试考生诚信档案》}：公开文件中明确写入该后果的情形是\textbf{弄虚作假、论文抄袭、虚报获奖或科研成果等学术不端行为}，不是单纯放弃\nref{S0527}\nref{S0553}。
\end{enumerate}

\textbf{真正的代价}：安徽师范大学2027年方案写明“推免工作结束后，若有学生放弃推免资格的，将等额扣除该生所在学院次年推免指标”\nref{S0646}——学校会追责到学院，学院会追责到你。
\nref{S0454}
\nref{S0567}
\nref{S0604}
\nref{S0639}
\nref{S0646}
\nref{S0527}
\nref{S0553}
\end{qa}

\iin{教育部学生服务与素质发展中心（学信网/研招网）}
\ith{违约}
\ist{资格与排名}
\begin{qa}{M0118}{推免资格名单上报教育部后还能改吗}
\textbf{上报备案后基本不能再改。} 要改必须在公示和备案截止前完成。

\textbf{全国层面}：教育部办公厅教学厅〔2014〕5号规定，各推荐高校应在推荐阶段通过“推免服务系统”将推荐办法及按名额遴选并公示的推免生名单报省级教育招生考试管理机构政策审核并向教育部备案，\textbf{“备案截止后，不再进行补充备案。最终推免生名单以'推免服务系统'备案信息为准，未经推荐高校公示及'推免服务系统'备案的推免生无效”}\nref{S0205}。

\textbf{校内程序层面}：云南师范大学2027年通知把时间卡得很紧——校内报名与资格审核9月3日前完成，学院推荐考核并公示不少于5日，学校统一公示不少于7日，“公示期间不得修改公示内容，如需变动，须另行公示7日，并对变动内容作出说明”，随后上报教育部推免服务系统\nref{S0568}。教育部规范管理文件同样要求“公示期间公示内容不得随意变动，如有变动，须对变动部分做出说明，并对变动内容另行公示”\nref{S0527}。

\textbf{接收端的连带条款}：河北大学2027年办法明确“未在教育部'推免服务系统'获得推免资格备案的考生，拟录取资格及其公示名单自动失效”\nref{S0604}；甘肃农业大学2027年管理办法也规定名单经公示无异议后通过推免服务系统上报教育部\nref{S0059}。

结论：唯一合法路径是“重新公示”，且必须在备案截止前。
\nref{S0205}
\nref{S0568}
\nref{S0527}
\nref{S0604}
\nref{S0059}
\end{qa}

\iin{中华人民共和国教育部}
\ith{上报流程}
\ist{资格与排名}
\begin{qa}{M0121}{学院说综合素质测评前X\%才有推免资格，综测怎么算}
\textbf{“综测”是学院自定的一套综合评价公式，没有全国统一算法——但骨架是固定的：学分绩×权重＋综合素质评价×权重。}

\textbf{骨架由校级规定，参数由学院定}。南开大学2021年实施办法第九条：“综合成绩＝计算学分绩×权重＋综合素质评价成绩×权重”；第十条：“\textbf{计算学分绩所包含的课程范围及计算学分绩在综合成绩中所占权重由各学院根据本学院情况自主设定，计算学分绩权重不得低于50\%}”；第十一条：综合素质评价指标体系由各学院设定，可将特殊学术专长、科研训练、参军入伍、志愿服务、国际组织实习等纳入\nref{S0561}。

\textbf{各校实际算法差异极大}：

\par\vspace{2pt}\nobreak\begin{center}\footnotesize
\begin{tabularx}{\textwidth}{>{\raggedright\arraybackslash}X>{\raggedright\arraybackslash}X}
\toprule
\textbf{学校（依据年份）} & \textbf{综合成绩结构} \\ \midrule
大连理工大学商学院 & T＝学习成绩×0.85＋思想素养（满分5分）＋综合素质和能力测评×0.1；思想素养 S＝(n1×3)/9＋思想品德评价成绩，n1为前三学年精神文明/社会工作/社会实践奖学金数量之和，评价分中班级民主评议1.2、辅导员评议0.6、班主任评议0.2\nref{S0549} \\
成都大学商学院 & 综合测评成绩＝学业成绩×90\%＋综合能力×10\%，总分100分；成绩相同则按学业成绩排序\nref{S0079} \\
温州医科大学（2027） & “学业成绩在推免综合成绩中的权重不低于90\%”；同一方面多项加分原则上只取一项\nref{S0563} \\
\bottomrule
\end{tabularx}
\end{center}\vspace{2pt}\par

\textbf{“综测前X\%”怎么看}：山东中医药大学中医学院2027年细则写成“\textbf{前4学年综合测评总成绩须排在本专业年级的前30\%}”\nref{S0639}；兰州大学按班型划线——国家级基地班前80\%、省级基地班前60\%、专业班前50\%\nref{S0345}。

\textbf{动作}：查清本院细则的综测公式、\textbf{学分绩权重}、以及综合素质评价由哪些子项构成。学业权重越低，加分项的杠杆越大（南开最低50\%\nref{S0561}，温州医科大学不低于90\%\nref{S0563}）。
\nref{S0561}
\nref{S0549}
\nref{S0079}
\nref{S0563}
\nref{S0639}
\nref{S0345}
\end{qa}

\iin{南开大学}
\ith{综测}
\ist{资格与排名}
\begin{qa}{M0122}{保研政策里德智体美劳具体卡哪些硬指标}
\textbf{“德”是硬门槛（一票否决），“智”是核心排序依据，“体”在不少学校也是硬条件，“美/劳”多为加分项而非门槛。}

\textbf{全国层面原文}：教育部关于规范和加强推免工作的通知要求“坚持\textbf{德智体美劳全面衡量，以德为先，把学生思想品德考核作为推免生遴选的重要内容和录取的重要依据}……\textbf{思想品德考核不合格者不予推荐和录取}”；同时“\textbf{将本科阶段学业综合成绩作为推免工作最基础的遴选指标}”；“可将学生参军入伍服兵役、参加志愿服务、到国际组织实习、科研成果、竞赛获奖等符合全面发展价值导向等因素纳入学校推免生遴选指标体系”\nref{S0318}\nref{S0527}。

\textbf{落到各校条款上}：

\par\vspace{2pt}\nobreak\begin{center}\footnotesize
\begin{tabularx}{\textwidth}{>{\raggedright\arraybackslash}X>{\raggedright\arraybackslash}X}
\toprule
\textbf{维度} & \textbf{典型条款} \\ \midrule
德 & 中南大学材料学院：“坚持德智体美劳全面衡量，以德为先……\textbf{对思想品德考核不合格者不予推荐}”\nref{S0273}；北京大学2027年办法：“把\textbf{思想政治素质和品德考核结果作为遴选推荐的首要依据}”\nref{S0553} \\
智 & 各校“学业综合成绩”排序；南开规定计算学分绩权重不得低于50\%\nref{S0561}；温州医科大学规定学业成绩权重不低于90\%\nref{S0563} \\
体 & 北京科技大学：“身心健康，\textbf{体育课成绩合格}”\nref{S0233}；中山大学强基转段：“国家规定的\textbf{学生体质健康测试达标（免测除外）}”\nref{S0248}；安徽师范大学把“体质健康测试成绩”列入复核与公示内容\nref{S0646} \\
美/劳 & 一般折算为加分而非门槛：中南大学材料学院设“文体素质比赛获奖加分”\nref{S0273} \\
\bottomrule
\end{tabularx}
\end{center}\vspace{2pt}\par

\textbf{结论}：能一票否决的是德（思想品德考核）和部分学校的体（体测达标/体育课合格）。智决定排序。美、劳在公开文件中基本以加分项形式出现，未见“美育不达标即取消资格”的条款。

\textbf{动作}：把本院细则的“必备条件”一节抄下来，标出哪些是“不予推荐”（一票否决），哪些是“加分/优先”。这两类条款后果完全不同。
\nref{S0318}
\nref{S0527}
\nref{S0273}
\nref{S0553}
\nref{S0563}
\nref{S0233}
\nref{S0248}
\nref{S0646}
\nref{S0040}
\nref{S0561}
\end{qa}

\iin{shpg（高校官网）}
\ith{综测}
\ist{资格与排名}
\section{名额分配}
\begin{qa}{M0061}{推免资格到底是学校评还是学院评？为什么学院之间名额差别那么大}
\textbf{两级分工：学校定规则和名额，学院定细则并组织遴选。} 资格最终由学校审定，但“谁能上”由学院的排序决定。

以广东外语外贸大学2026年实施办法为例\nref{S0134}：
\begin{itemize}
  \item 学校推免生遴选工作领导小组负责“推免名额分配方案、推免名单等重要事项的研究审定”；
  \item 领导小组办公室（设在教务部）负责“审核各学院推免工作细则、下达各学院推免名额、监督和检查各学院推荐遴选工作、审核并上报学校推免生数据”；
  \item 各学院推免生遴选工作小组“负责制定本单位推免工作实施细则，明确推免名额分配原则、推免生评价指标体系等”，并“审议通过并提交本单位拟推免名单及推荐次序”。
\end{itemize}

南开大学2021年实施办法的结构相同：校级定实施办法与名额分配方案，学院制定《推免细则》并报校小组审核后公示\nref{S0561}。内蒙古大学的写法是“学校制定年度推免生名额分配办法，确定学院推免生名额”，再由“学院统筹使用学校分配的推免生名额，并制定学院推免生名额分配办法报学校备案后执行”\nref{S0544}。

\textbf{学院之间名额差别大的原因}：教育部给学校是总额比例——教学〔2006〕14号第十二条规定“教育部批准设立研究生院的高等学校一般按应届本科毕业生数的15\%左右确定；未设立研究生院的'211工程'建设高校一般按5\%左右；其他高校一般按2\%确定”，基地班“按教育部批准的基地班招生人数的50\%左右单独增加”\nref{S0454}。学校再往下分时考虑的因素各校不同：广外看“应届本科毕业生人数、硕士点、本科专业数量以及上年度推免生名额使用情况”\nref{S0134}；内蒙古大学看“人才培养基地、一流专业、基础学科及国家战略和民生领域急需的相关学科专业建设水平，以及人才培养质量、应届毕业生人数”\nref{S0544}；云南师范大学按“卓越人才培养计划”“直博生”“专业建设奖励”三类切块，其中专业建设奖励“根据国家和区域发展对专业人才的需求、专业建设水平、升学率等情况分配”\nref{S0568}。
\nref{S0134}
\nref{S0561}
\nref{S0544}
\nref{S0454}
\nref{S0568}
\nref{S0273}
\end{qa}

\iin{广东外语外贸大学}
\ith{资格评定}
\ist{资格与排名}
\begin{qa}{M0062}{我们学院保研名额是怎么分的？按专业人数还是按往年升学率}
\textbf{名额是学校下达到学院、学院再分到专业的，切分规则各院自定，公开程度不一。}

学院内部怎么切，通常是这几类依据的加权：当年应届毕业班人数、一流专业/硕士点/重点学科建设需要、往年使用情况。成都大学商学院2027届公布了一份可直接对照的分配表\nref{S0079}：

\par\vspace{2pt}\nobreak\begin{center}\footnotesize
\begin{tabularx}{\textwidth}{>{\raggedright\arraybackslash}X>{\raggedright\arraybackslash}X>{\raggedright\arraybackslash}X}
\toprule
\textbf{专业} & \textbf{推免指标} & \textbf{候补指标} \\ \midrule
会计学 & 3 & 2 \\
审计学 & 2 & 1 \\
财务管理 & 2 & 1 \\
工商管理 & 3 & 2 \\
国际经济与贸易 & 1 & 1 \\
金融科技 & 2 & 1 \\
\bottomrule
\end{tabularx}
\end{center}\vspace{2pt}\par

它写明的依据是“结合学院各本科专业应届毕业班人数、本科一流专业建设及所对应的硕士学位点建设、重点学科建设需要”，并规定各专业名额之和不得超过学校下达总数，候补人数按各专业推免名额数的50\%计算\nref{S0079}。

其他写法：安徽师范大学2027年方案把普通类名额拆成“基础指标”和“综合指标”，\textbf{基础指标＝（教育部下达的普通类名额－综合指标名额）÷ 全校应届毕业学生总数 × 学院人数}，同时“综合考察学院应届本科预计毕业生人数、人才培养质量、专业建设水平、学科发展及社会评价等因素”分配综合指标\nref{S0646}。云南师范大学2027年通知把名额分成卓越人才培养计划156个、直博生8个、专业建设奖励226个三类，其中专业建设奖励按“专业建设水平、升学率”分配到各单位\nref{S0568}。温州医科大学则按“优先向一流学科建设、国家战略紧缺医学专业、拔尖创新人才倾斜”，学院层面按1:1.2确定推荐候选人名额\nref{S0563}。

\textbf{结论}：没有“按人数还是按升学率”的统一答案，多数学校是“人数打底＋学科/专业建设加权”。你对分配有疑问，应在学院细则公示期提出。
\nref{S0079}
\nref{S0646}
\nref{S0568}
\nref{S0563}
\nref{S0134}
\end{qa}

\iin{成都大学商学院}
\ith{名额分配}
\ist{资格与排名}
\begin{qa}{M0071}{保研加分项一般都有哪些？哪些是通用的哪些是学校自定}
\textbf{通用的是“五类加分方向”，各校自定的是目录、分值和上限。}

\textbf{全国层面的指引}：教育部关于规范和加强推免工作的通知要求，推荐高校“可将学生参军入伍服兵役、参加志愿服务、到国际组织实习、科研成果、竞赛获奖等符合全面发展价值导向等因素纳入学校推免生遴选指标体系”，同时“要合理设置各遴选指标所占权重及单项指标上限分值，学生在某一方面中有多项加分情况时，原则上只取一项，不得为仅符合单一或部分遴选指标的学生单列计划或破格推荐”\nref{S0318}\nref{S0527}。所以\textbf{通用方向就这五类}，其余全凭学校自定。

\textbf{各校自定的真实样貌}：

\par\vspace{2pt}\nobreak\begin{center}\footnotesize
\begin{tabularx}{\textwidth}{>{\raggedright\arraybackslash}X>{\raggedright\arraybackslash}X>{\raggedright\arraybackslash}X}
\toprule
\textbf{学校（依据年份）} & \textbf{加分结构} & \textbf{上限} \\ \midrule
中南大学材料学院（2027细则） & 学科竞赛、学术科研成果、荣誉称号、文体素质比赛 & 学院评价加分累计≤2分；学校专项加分各类型累计≤10分（不含退役复学）\nref{S0273} \\
湖南大学（2027加分公示） & 创新创业赛事（挑战杯）、体育类、艺术类国家级及以上奖励、志愿服务 & 挑战杯国家级特等奖排名第1加6分、三等奖排名第1加2分\nref{S0541} \\
内蒙古大学（实施办法） & 核心期刊论文、全国赛获奖、参军入伍、国际组织实习、志愿服务、体育比赛 & 论文≤0.15、全国赛一等≤0.15/二等≤0.12/三等≤0.1\nref{S0544} \\
安徽师范大学（2027方案） & 学术专长、入伍建功、国际交流、志愿服务 & 入伍荣立三等功及以上专长分值300分，国际组织实习3个月及以上100分\nref{S0646} \\
\bottomrule
\end{tabularx}
\end{center}\vspace{2pt}\par

共同点：同一类别多项\textbf{只取最高、不累计}；科研论文和竞赛要经专家审核小组鉴定，未通过的不计入。
\nref{S0318}
\nref{S0527}
\nref{S0273}
\nref{S0541}
\nref{S0544}
\nref{S0646}
\end{qa}

\iin{shpg（高校官网）}
\ith{加分项}
\ist{资格与排名}
\begin{qa}{M0082}{保研率是按专业算还是按学院算？转专业后算哪个专业的名额}
\textbf{名额按学院切、排名按专业排；转专业后一般随现专业。}

\begin{itemize}
  \item \textbf{名额}：学校把名额下达到学院，学院再分到专业。内蒙古大学明确“学校确定学院推免生名额。学院统筹使用学校分配的推免生名额”，同时“各学院不得将推免生名额跨学院使用”\nref{S0544}。所以跨学院的名额不能互借。
  \item \textbf{排名}：多数学校按“专业”排。中南大学材料学院规定推免生成绩“按专业进行排名”\nref{S0273}；成都大学商学院写的是“本专业推荐对象基数的前20\%”\nref{S0079}；南开大学写成“本专业本年级具有我校学籍学生的前80\%以内（以排序时在校正式学籍学生数为基准）”\nref{S0561}。
  \item \textbf{所以“保研率”是学院的指标、专业的基数}：同一个学院下不同专业的保研率可能差好几倍。
  \item \textbf{转专业后算哪个专业}：主流是随\textbf{现在所在专业}。南京工业大学2024年通知在推荐对象里明确“含校内转专业学生”\nref{S0527}，即转专业学生照常参与现专业推免。但“原来的课算不算”是另一个问题——大连理工大学商学院的规定是“转专业学生按目前所学专业对应的必修课程范围计算，若所修课程不在规定课程范围内，则该课程不参与计算”\nref{S0549}。
\end{itemize}

结论：名额归属看学院文件的二次分配表，排名归属看“现专业＋现年级＋在籍人数”。转专业学生一定要书面问清“我按哪个专业、哪一届排名”。
\nref{S0544}
\nref{S0273}
\nref{S0079}
\nref{S0561}
\nref{S0527}
\nref{S0549}
\end{qa}

\iin{内蒙古大学}
\ith{保研率}
\ist{资格与排名}
\begin{qa}{M0083}{大类招生的，专业分流后才定排名，保研名额怎么分}
\textbf{现行公开文件中未见针对“大类招生分流后名额归属”的全国统一规定}，须以本校当年文件为准。

可以确定的通用原则有三条，你可以据此向学院提问：

\begin{enumerate}
  \item \textbf{排名基数用“现专业在籍学生数”}。成都大学商学院写的是“本专业推荐对象基数”\nref{S0079}；南开大学写的是“本专业本年级具有我校学籍学生的前80\%以内（以排序时在校正式学籍学生数为基准）”\nref{S0561}——注意“以排序时”这个措辞，意味着基数按\textbf{排名时点}的实际在籍人数算，而不是入学时的招生人数。
  \item \textbf{名额是先到学院、再由学院分到专业}。安徽师范大学2027年方案的基础指标＝（普通类名额－综合指标名额）÷ 全校应届毕业生总数 × 学院人数，学院再“根据本院推免工作实施细则及名额分配方案，完成各专业推免名额分配”\nref{S0646}。所以大类分流后，学院要重新在分流后的各专业间切一次。
  \item \textbf{各学院须制定并公示实施细则}。广东外语外贸大学要求“学院推免生遴选工作小组负责制定本单位推免工作实施细则，明确推免名额分配原则、推免生评价指标体系等，报领导小组办公室备案审查”\nref{S0134}，分流规则应当写在里面。
\end{enumerate}

\textbf{实操建议}：直接书面问学院教学办三个问题——我按分流后的哪个专业排名？分母含不含降级、休学、转出学生？大类阶段（未分流时）的成绩怎么计入？这三个答案决定了你的位置。
\nref{S0079}
\nref{S0561}
\nref{S0646}
\nref{S0134}
\end{qa}

\iin{成都大学商学院}
\ith{大类分流}
\ist{资格与排名}
\begin{qa}{M0088}{基地班实验班的保研名额是单列的吗}
\textbf{部分班型确有单列或倾斜，但不是自动单列——最终“由学校统筹安排”。}

\textbf{全国层面的依据}：教育部《全国普通高等学校推荐优秀应届本科毕业生免试攻读硕士学位研究生工作管理办法（试行）》（教学〔2006〕14号）第十二条第（四）项明确：经教育部确定的人文、理科等人才培养基地的高等学校，“按教育部批准的基地班招生人数的50\%左右，单独增加推免生名额，由学校统筹安排”\nref{S0454}。这句话有两层意思：基地班的名额是\textbf{额外增加的}，但\textbf{怎么用由学校定}，不等于自动划归该班。

\textbf{各校落地时的实际差异}：

\par\vspace{2pt}\nobreak\begin{center}\footnotesize
\begin{tabularx}{\textwidth}{>{\raggedright\arraybackslash}X>{\raggedright\arraybackslash}X}
\toprule
\textbf{学校/班型} & \textbf{口径} \\ \midrule
中南大学材料学院 卓越人才培养实验班 & 学业成绩须“在本专业卓越班应届生中排名前60\%”，平行班和国际班则是前35\%\nref{S0273} \\
中国药科大学 理科基地班/生科基地班/拔尖计划 & 本科推荐环节成绩排名前30\%（国家级一流本科专业前10\%、其他专业前8\%）\nref{S0417}\nref{S0418} \\
兰州大学 & 国家级人才培养基地班前三年综合测评前80\%、省级基地班前60\%、专业班前50\%（五年制为前四年）\nref{S0345} \\
\bottomrule
\end{tabularx}
\end{center}\vspace{2pt}\par

可以看到：基地班/实验班的门槛值通常比普通班宽松，等于在同一学院内“单独划线”。这实质上是名额向这些班级倾斜的一种方式。

\textbf{结论}：基地班、实验班的名额是否单列，取决于本校是把名额直接下达到班级，还是只给宽松的排名门槛。前者才是真单列。请查本院细则里的名额分配表有没有按班单列。
\nref{S0454}
\nref{S0273}
\nref{S0417}
\nref{S0418}
\nref{S0345}
\nref{S0134}
\end{qa}

\iin{教育部学生服务与素质发展中心（学信网/研招网）}
\ith{特殊班级}
\ist{资格与排名}
\begin{qa}{M0090}{卓越班珠峰计划的保研率是不是更高}
\textbf{“珠峰计划”推免在现行公开文件中未见全国统一规定}，须以本校当年文件为准。至于“卓越”类班型，公开文件显示它们确实享有一些制度性倾斜，但倾斜方式不一，不能直接等同于“保研率更高”。

可以看到的四类倾斜：

\begin{enumerate}
  \item \textbf{选拔门槛放宽}。中南大学材料科学与工程学院2027年细则规定：平行班和国际班学生学业成绩须“在本专业应届生中排名前35\%”，而“材料科学与工程卓越人才培养实验班学生须在本专业卓越班应届生中排名前60\%”\nref{S0273}。
  \item \textbf{名额单独切块}。云南师范大学2027年通知把“卓越人才培养计划推免名额为156个”单列，明确“分配给实施'卓越教师''卓越工程师''卓越法治人才''卓越新闻传播人才'培养计划的各单位”\nref{S0568}。
  \item \textbf{推荐类型单列}。北京科技大学2027年通知列出推荐类型包括“普通类、人才培养改革类（理科试验班、工科试验班、黄昆英才班、钢铁脊梁计划、科史新星计划、青马人才工程、材料高精尖班、机器人科创班、卓越创新班、纳米试验班）、研究生支教团等”，各类型“有不同的推荐条件，见负责推荐单位的具体通知，学生申请其中一项推荐类型，参加相应类型的遴选”\nref{S0233}。
  \item \textbf{培养层次倾斜}（非推免率）。北京师范大学的“卓越博士”计划是入学后的奖学金计划，不是推免名额\nref{S0423}。
\end{enumerate}

\textbf{结论}：门槛更低、名额单列，都会提高实际获得推免的概率，但三者不是一回事。要判断“保研率是不是更高”，只能向学院索要该班近三年的\textbf{实际推免人数 ÷ 班级人数}。
\nref{S0273}
\nref{S0568}
\nref{S0233}
\nref{S0423}
\nref{S0345}
\end{qa}

\iin{中南大学}
\ith{特殊班级}
\ist{资格与排名}
\begin{qa}{M0096}{研究生支教团的保研名额是额外的还是占学院名额}
\textbf{是额外的，不占学院名额。} 这一点各校表述高度一致。

\begin{itemize}
  \item 郑州大学团委通知：拟公开招募27名具备推免条件的2026届毕业生，“参加研究生支教团的2026届本科毕业生不占用我校今年的免试推荐研究生名额，由教育部专项下拨”\nref{S0597}。
  \item 安徽工业大学2027年招募方案：“第29届研究生支教团招募指标由教育部专项下拨为4名。研究生支教团推免指标不占用学校计划推荐免试研究生的名额”\nref{S0589}。
  \item 西南交通大学第29届招募通知：“2026年招募计划为18人（该推免指标由教育部专项下拨，不占用学校下达给各学院的研究生免试推荐名额）”\nref{S0061}。
  \item 兰州大学第29届服务队通知：“我校第29届研究生支教团招募名额为65名（由教育部专项下拨，在读硕、博士研究生申报不占该名额）”\nref{S0345}。
\end{itemize}

\textbf{但要注意两点}：

\begin{enumerate}
  \item \textbf{名额来源可能是“部分专拨＋部分配套”}。云南师范大学2027年通知写明：“'研究生支教团'推免名额为15个，包括教育部下达的12个专项名额和学校配套的3个名额”\nref{S0568}——配套部分是不是占学校总名额，要问本校团委。
  \item \textbf{门槛不降低}。多校明确“研究生支教团推免生综合测评成绩底线应与普通类推免生综合测评成绩底线保持一致”\nref{S0370}\nref{S0345}\nref{S0061}；安徽工业大学的招募对象也要求“具备推荐本校免试攻读研究生条件”\nref{S0589}。也就是说，你要先够得上普通推免的门槛，才能去竞争这个额外名额。
\end{enumerate}

\textbf{组织单位}：由校团委负责招募选拔，与教务处/研究生院的普通推免流程分开走。
\nref{S0597}
\nref{S0589}
\nref{S0061}
\nref{S0345}
\nref{S0568}
\nref{S0370}
\end{qa}

\iin{郑州大学}
\ith{支教团}
\ist{资格与排名}
\begin{qa}{M0112}{有人放弃保研名额，名额会回到学院还是学校收回}
\textbf{分两段看：截止前退回学校统筹，截止后名额作废、不得转让。学院通常留不住。}

\textbf{全国层面的硬规定}：教育部教学〔2006〕14号第十三条——“高等学校不得将推免生名额跨学校使用……\textbf{截至教育部规定时间后，被确定的推免生如果无招生单位接收或本人放弃，则该生推免生资格作废，名额不得转让}”\nref{S0454}。

\textbf{校内层面的两种做法}：
\begin{itemize}
  \item \textbf{退回学校重新分配}。内蒙古大学实施办法第八条写得很清楚：“在教育部规定的推免工作截止前，被确定的推免生如果无招生单位接收或本人放弃，该推免生名额退回\textbf{学校}，由学校推免工作领导小组重新分配。截止教育部规定时间后……则该生推免资格作废”\nref{S0544}。云南师范大学2027年通知同样规定“各单位分配名额未使用完的，由学校收回，并经学校推免工作领导小组审定后，面向学校'冲A专业'以及符合国家和地方重大发展需求的专业重新分配”\nref{S0568}。
  \item \textbf{校内顺次递补}。公示期内的放弃，通常由本校候补按综合成绩依次补上。南京林业大学风景园林学院2027年公示明确“若已获推荐资格者因故放弃资格，或在公示期内被查实不符合推荐条件，其空余名额将依据综合成绩排序，由后续候选人依次递补”\nref{S0422}。
\end{itemize}

\textbf{还有一笔账}：安徽师范大学2027年方案规定“推免工作结束后，若有学生放弃推免资格的，将等额扣除该生所在学院次年推免指标”\nref{S0646}。

结论：放弃后的名额流向是“先校内递补、截止后收归学校、最终截止后作废”，\textbf{不会默认留在学院}。
\nref{S0454}
\nref{S0544}
\nref{S0568}
\nref{S0422}
\nref{S0646}
\nref{S0561}
\end{qa}

\iin{教育部学生服务与素质发展中心（学信网/研招网）}
\ith{名额分配}
\ist{资格与排名}
\begin{qa}{M0113}{学校把名额调给别的学院，原学院的学生能反对吗}
\textbf{名额在校内怎么切分，属于学校自主权，学生一般不能“反对”；但你可以在程序层面申诉。}

\textbf{为什么不能反对}：推免名额是教育部下达到学校、再由学校向下分配的。内蒙古大学实施办法规定“学校……制定年度推免生名额分配办法，确定学院推免生名额”，学院“统筹使用学校分配的推免生名额”，并且“各学院不得将推免生名额跨学院使用”\nref{S0544}。云南师范大学2027年通知进一步写明“各单位分配名额未使用完的，由学校收回，并经学校推免工作领导小组审定后，面向学校'冲A专业'以及符合国家和地方重大发展需求的专业重新分配”\nref{S0568}。这都说明名额的再分配权在学校，不在学院，也不在学生。

\textbf{你能做的是走程序}：
\begin{itemize}
  \item \textbf{学院复议}。湖南大学2027年工作安排规定：如学生对推荐结果有异议，填写《湖南大学推免生复议申请表》（一式三份），向所在单位推荐工作小组申请复议；工作小组审定并在表中详细记录复议过程及结论，组长签字并加盖公章，\textbf{“复议结论应在推荐名单报送推免生遴选工作领导小组之前告知申请人”}\nref{S0540}\nref{S0537}。
  \item \textbf{校级申诉渠道}。北京大学2027年推免办法要求“各院系应畅通学生咨询及申诉渠道，对学生提出的对推免相关的意见或申诉应及时调查处理并予以反馈”\nref{S0553}；云南师范大学要求各学院组建专门的推免工作监督小组“受理推免工作中的异议和申诉”\nref{S0568}。
  \item \textbf{通用学生申诉}。教育部令第41号第五十九条至第六十一条：学校应成立学生申诉处理委员会；学生对处理或处分决定有异议的，可在接到决定书之日起10日内提出书面申诉，申诉处理委员会应在接到书面申诉之日起15日内作出复查结论，情况复杂可延长15日\nref{S0040}。
\end{itemize}

\textbf{结论}：名额调整本身你难以推翻，但如果调整导致\textbf{已公示结果变更而未重新公示}，或程序违反本校文件，你可以据此申诉。
\nref{S0544}
\nref{S0568}
\nref{S0540}
\nref{S0537}
\nref{S0553}
\nref{S0040}
\end{qa}

\iin{内蒙古大学}
\ith{名额调整}
\ist{资格与排名}
\begin{qa}{M0115}{港澳台侨学生的推免名额是单列的吗}
\textbf{未见“港澳台侨学生单列推免名额”的公开规定；相反，多数学校的推免对象条款把他们明确纳入了。}

\begin{itemize}
  \item 教育部教学〔2006〕14号第十四条对推荐对象的表述是“纳入国家普通本科招生计划录取的应届毕业生”\nref{S0454}。
  \item 各校据此写的推荐对象条款，有的直接点了港澳台侨：成都大学商学院“推荐对象为本校通过普通高考（含保送方式）招生录取\textbf{或面向港澳台华侨学生招生录取}的全日制应届本科毕业生”\nref{S0079}；山东中医药大学中医学院“选拔对象为本校通过普通高考（含保送方式）招生录取\textbf{或面向港澳台华侨学生招生录取}的全日制本科生”\nref{S0639}；武汉轻工大学2027年办法的申请条件同样写作“通过普通高考（含保送方式）招生录取\textbf{或面向港澳台华侨学生招生录取}的全日制本科生”\nref{S0409}。
\end{itemize}

也就是说，这类学生属于\textbf{普通推免序列}，与内地（大陆）学生使用同一套名额和同一套排名口径，没有单独的“港澳台名额”。

\textbf{但有两处需要你自己核实}：
\begin{enumerate}
  \item \textbf{身份对“定向就业/骨干计划”的影响}。少数民族高层次骨干人才计划等专项对生源地有明确要求\nref{S0573}，港澳台侨身份通常不适用。
  \item \textbf{个别学校的推荐对象条款可能另有排除}。广东外语外贸大学排除的是“辅修学位、专升本、第二学士学位学生、本硕项目的国际班学生和独立学院学生”，未提港澳台侨\nref{S0134}。
\end{enumerate}

结论：一般纳入普通序列统一排名，无单列名额。以本校当年文件的“推荐对象”条款为准。
\nref{S0079}
\nref{S0639}
\nref{S0409}
\nref{S0454}
\nref{S0134}
\nref{S0573}
\end{qa}

\iin{成都大学商学院}
\ith{港澳台}
\ist{资格与排名}
\section{成绩折算}
\begin{qa}{M0067}{四六级成绩算不算保研加分或门槛}
\textbf{主要是“门槛”，不是“加分”。} 四六级在推免里最常见的用法是设一条及格线，而不是把分数折算成加分；把六级高分直接折算加分的做法在公开文件中较少见。

\textbf{各校门槛写法对照}：

\par\vspace{2pt}\nobreak\begin{center}\footnotesize
\begin{tabularx}{\textwidth}{>{\raggedright\arraybackslash}X>{\raggedright\arraybackslash}X}
\toprule
\textbf{学校（依据年份）} & \textbf{外语门槛} \\ \midrule
内蒙古大学 & 非英语专业四级≥425分；艺术类、体育类及高考英语未按100\%计入总分者355分或英语应用能力考试50分以上；六级≥425分者同等条件下优先\nref{S0544} \\
广东外语外贸大学（2026办法） & 非外语类（除国际商务、外交学、教育学、艺术类）六级≥425分，或托福≥90，或雅思≥6.5；其他外语语种和艺术类四级≥425\nref{S0134} \\
成都大学商学院（2027届） & 非英语语种非语言类大学外语平均≥75分；其他本科学生四级≥425分，或雅思5.5，或托福66\nref{S0079} \\
燕山大学（2027通知） & 大学英语四级不低于425分（小语种取得相应语种等级考试合格证书）\nref{S0533} \\
温州医科大学（2027通知） & 全国大学英语六级≥425分；雅思6.0或托福75以上视为符合\nref{S0563} \\
兰州大学 & 非外语类四级或六级达425分，其他语种成绩达满分值60\%\nref{S0345} \\
\bottomrule
\end{tabularx}
\end{center}\vspace{2pt}\par

\textbf{加分方面}：目前能确认的公开写法只有“同等条件下优先”（内蒙古大学对六级≥425分者\nref{S0544}），没有见到把四六级分数计入加分目录的普遍做法。

\textbf{注意}：接收院校（夏令营/预推免）对外语另有要求，通常比本校门槛更严。例如北京大学多个院系要求四级525分或六级450分以上\nref{S0213}。
\nref{S0544}
\nref{S0134}
\nref{S0079}
\nref{S0533}
\nref{S0563}
\nref{S0345}
\nref{S0213}
\end{qa}

\iin{内蒙古大学}
\ith{英语门槛}
\ist{资格与排名}
\begin{qa}{M0074}{学生会主席有没有保研加分}
\textbf{“学生会主席”这个身份本身在多数学校不加分；能加分的是与之相关的“荣誉称号”，且分值通常很小。}

\textbf{全国层面的约束}：教育部推免规范管理文件要求“合理设置各遴选指标所占权重及单项指标上限分值，学生在某一方面中有多项加分情况时，原则上只取一项，\textbf{不得为仅符合单一或部分遴选指标的学生单列计划或破格推荐}”\nref{S0318}\nref{S0527}。这几句话实际上限制了“靠学生干部身份拿大分”的空间。

\textbf{各校实际写法}：

\begin{itemize}
  \item \textbf{按荣誉称号加分，不按职务加分}。中南大学材料科学与工程学院2027年加分细则里，“荣誉称号加分”包含：国家级荣誉称号加2.0分；省部级如省优秀共青团员、省优秀共青团干部加0.8分；校级中“优秀学生干部标兵”“优秀团干标兵”加0.4分，“优秀学生、优秀学生干部、优秀团员、优秀团干、优秀学生会干部”加0.3分，院级优秀学生干部加0.1分；\textbf{“有多项荣誉获得者，不累计加分，取最高分”}\nref{S0273}。
  \item \textbf{折算成“思想素养/素质”评分}。大连理工大学商学院把“思想素养”满分设为5分，算法为 S=(n1×3)/9+思想品德评价成绩，其中 n1 是前三学年所获精神文明奖学金、社会工作奖学金、社会实践奖学金数量之和；思想品德评价成绩由班级民主评议（满分1.2）、辅导员评议（0.6）、班主任评议（0.2）构成\nref{S0549}。
  \item \textbf{志愿服务/社会工作可单列加分项}。中南大学的专项评价成绩明确包含“参与社会工作志愿服务加分”，由学生工作部认定公示\nref{S0273}。
\end{itemize}

\textbf{结论}：当学生会主席不会直接加分，但如果这个岗位带来了校级“优秀学生干部”荣誉，很多学校能折成0.1—0.4分级别的加分——在名次密集处这仍有意义，但分值上限普遍很低。
\nref{S0273}
\nref{S0549}
\nref{S0318}
\nref{S0527}
\end{qa}

\iin{中南大学}
\ith{社工加分}
\ist{资格与排名}
\begin{qa}{M0075}{志愿服务时长算不算保研加分}
\textbf{看本校：有的学校把志愿服务写进加分目录，有的只是“同等条件优先”，还有的明确不认时长、只认表彰。}

\textbf{全国层面的指引}：教育部推免规范管理文件把“参加志愿服务”列为可纳入遴选指标体系的五类因素之一，但同时要求“合理设置各遴选指标所占权重及单项指标上限分值”，且同一方面多项只取一项\nref{S0318}\nref{S0527}。所以能做，但怎么算由学校定。

\textbf{三种真实写法}：

\begin{enumerate}
  \item \textbf{列为专长加分项}。中南大学材料学院2027年细则的专项评价成绩里包含“参与社会工作志愿服务加分”，由学生工作部负责认定并公示\nref{S0273}。湖南大学2027年单独发布了包含“志愿服务加分”的公示名单\nref{S0541}。
  \item \textbf{只认表彰、不认时长}。安徽师范大学2027年方案规定：志愿服务类要“入学以来每学年坚持参加志愿服务活动\textbf{并获得省级以上志愿服务先进个人表彰}，或在\textbf{国家级重大活动中积极参加志愿服务表现突出}”，专长对应分值40分\nref{S0646}。
  \item \textbf{作为同等条件优先}。兰州大学第29届支教团通知把“积极参加志愿服务、兰州大学'四星级'及以上志愿者”列为同等条件下优先考虑的情形\nref{S0345}。
\end{enumerate}

\textbf{结论}：单纯累计志愿时长在多数学校不直接折算为推免加分；能被认定的是“有组织的志愿服务项目＋表彰/星级/项目认定”。

\textbf{注意}：“有组织”三个字很关键——北京大学2027年推免办法把可加分的情形写作“参加院系认可的、\textbf{有组织的}志愿服务”\nref{S0553}。自行找的志愿活动可能不被认定。
\nref{S0273}
\nref{S0541}
\nref{S0646}
\nref{S0345}
\nref{S0553}
\nref{S0318}
\end{qa}

\iin{中南大学}
\ith{志愿加分}
\ist{资格与排名}
\begin{qa}{M0091}{定向生委培生能不能保研}
\textbf{一般不能；除非原定向单位书面同意，或本校文件明确允许。}

\textbf{制度根源}：教育部教学〔2006〕14号第十四条要求推免对象为“纳入国家普通本科招生计划录取的应届毕业生”，而定向生、委培生入学时已与定向单位签订协议，其升学受协议约束\nref{S0454}。

\textbf{各校的实际处理}：

\begin{itemize}
  \item \textbf{明确不接收}。中山大学2027年接收推免生报名通知写明：“我校不接收未按国家有关规定完成履约责任（含违约）的'农村订单定向生''教育部直属师范院校公费师范生'等推免生申请，一经查实，将取消其申请资格、录取资格、入学资格，已入学的取消学籍。”\nref{S0166}
  \item \textbf{留了口子}。浙江师范大学2027年推免工作通知的推荐对象为“2023级（四年制）或2022级（五年制）全日制普通本科毕业生（\textbf{不含定向生，但原定向单位书面同意推荐或硕师计划有关教育局指定的除外}）”\nref{S0569}。也就是“定向单位书面同意”可以解套。
  \item \textbf{师范类另有通道}。云南师范大学2027年通知把“公费师范生、优师计划师范生”排除在普通推免对象之外\nref{S0568}，但公费师范生有独立的“本研衔接公费教育”路径（见免费医学生/公费农科生条目）。
\end{itemize}

\textbf{结论}：定向生、委培生能否推免，取决于\textbf{两件事}：定向协议里对升学的约定，以及本校推荐对象条款有没有写“原定向单位书面同意可除外”。先看协议，再问教务处，最后拿到定向单位的书面意见。
\nref{S0454}
\nref{S0166}
\nref{S0569}
\nref{S0568}
\end{qa}

\iin{教育部学生服务与素质发展中心（学信网/研招网）}
\ith{定向生}
\ist{资格与排名}
\begin{qa}{M0092}{免费医学生公费农科生能保研吗}
\textbf{免费医学生、公费农科生属于定向培养，普通推免通道基本走不通；能免试读研的是“本研衔接/定向专项”这条独立路径。}

\textbf{一、普通推免受限于履约。} 中山大学2027年接收通知明确不接收“未按国家有关规定完成履约责任（含违约）的'农村订单定向生''教育部直属师范院校公费师范生'”\nref{S0166}。电子科技大学2026年招生章程对统考口径写明：“服务期限内的农村订单定向免费培养医学毕业生只能报考临床医学专业学位硕士全科医学方向，且报考类别为定向就业；考生在报考时必须提供定向就业县（市、区）卫生健康行政部门出具的\textbf{同意报考……书面同意函}”\nref{S0493}。

\textbf{二、公费师范生有专门通道，但不是普通推免。} 国务院办公厅转发的《教育部直属师范大学本研衔接师范生公费教育实施办法》（国办发〔2024〕27号）：从2024年起在六所部属师范大学实施本研衔接公费教育，支持符合条件的公费师范生\textbf{免试攻读本校全日制教育硕士研究生}再履约任教\nref{S0456}\nref{S0355}。教育部办公厅2025—2029年就业工作通知明确：办法印发前已签约的本科层次公费师范生，本科三年级结束后可申请，“\textbf{经学校对其进行推免综合考核，取得学校推免资格并达到学校硕士研究生录取标准}”，报生源所在省份省级教育行政部门同意并签署补充协议\nref{S0044}。各校据此制定专门办法，如华中师范大学2027年公费师范生免试攻读教育硕士办法\nref{S0559}、东北师范大学2027年办法（要求“毕业后按协议履约任教满一学期，经任教学校考核合格并批准报考”\nref{S0609}）。

\textbf{三、公费农科生}：现行公开文件中未见对应的全国性推免或本研衔接规定，须以本省和本校当年文件为准。

\textbf{结论}：免费医学生看定向卫健部门是否书面同意；公费师范生走“本研衔接”专项（仍需取得推免资格）；公费农科生暂无公开统一路径。
\nref{S0166}
\nref{S0493}
\nref{S0456}
\nref{S0355}
\nref{S0044}
\nref{S0559}
\nref{S0609}
\end{qa}

\iin{中山大学}
\ith{公费生}
\ist{资格与排名}
\begin{qa}{M0093}{国防生现在还能保研吗}
\textbf{2017年起国家已不再招收和选拔国防生，所以“国防生保研”这条路本身已经不存在。}

这一变化有明确的政策来源。中国政府网转载的2017年军队院校招生政策解读中，招生办负责人答记者问时说明：“根据教育部、军委政治工作部、军委训练管理部通知，\textbf{从2017年起不再从普通高中毕业生中定向招收国防生，也不再从在校生中直接选拔国防生}”\nref{S0572}。此后在校本科生中不再新增国防生身份，原有的国防生政策也随培养周期结束而自然终止。

\textbf{现在与军队相关的读研通道是另外三种}：

\begin{enumerate}
  \item \textbf{退役大学生士兵专项计划}。教育部每年实施，招生计划“在全国硕士研究生招生计划总规模内\textbf{单列下达、专项使用，不得调整为普通计划或其他专项计划}”，招生单位自主确定该计划考生的初试成绩要求（一般不高于本单位相应学科专业普通计划考生初试成绩要求）\nref{S0052}\nref{S0440}。这是\textbf{统考}通道，不是推免。
  \item \textbf{免初试攻读硕士研究生}。《2027年全国硕士研究生招生工作管理规定》第二十五条：服役期间获得三等战功、二等功以上奖励或者二级以上表彰，符合报考条件的退役人员，“\textbf{可申请免初试攻读硕士研究生}”，须在教育部规定的统考报名时间内登录推免服务系统报名，招生单位组织复试，复试合格者予以拟录取\nref{S0442}。
  \item \textbf{军队院校接收地方院校推免生}。如国防科技大学接收地方院校推荐免试硕士研究生（含直博生、\textbf{参军入伍生}），设有“国防科工计划”，面向地方院校推免生\nref{S0420}。这是“参军入伍生”路径，与旧“国防生”不同。
\end{enumerate}

\textbf{结论}：如果你是原国防生，按当年的协议和政策执行；如果你是现在在读的本科生，对应的是“大学生应征入伍→退役大学生士兵计划”或“免初试”通道。
\nref{S0572}
\nref{S0052}
\nref{S0440}
\nref{S0442}
\nref{S0420}
\end{qa}

\iin{中国政府网}
\ith{国防生}
\ist{资格与排名}
\begin{qa}{M0100}{受过违纪处分还能保研吗}
\textbf{分两种写法，结果截然相反；而且拿到资格后再受处分一样会被取消。}

\textbf{一、校内基本条件的两种写法}

\begin{itemize}
  \item \textbf{“不在处分期内”}——解除后仍可申请。燕山大学2027年通知的遴选条件第5条是“品行表现优良，\textbf{不在违纪处分期内}”\nref{S0533}。中山大学2026年强基计划转段方案写得更细：“在校期间虽受到处分（因考试作弊或学术不端行为受到处分除外），但在申请转段资格前\textbf{已依法依规解除}”\nref{S0248}——即一般处分解除后可参与，考试作弊或学术不端除外。
  \item \textbf{“无任何违法违纪受处分记录”}——一刀切。武汉轻工大学2027年办法要求“品行表现优良，无任何违法违纪受处分记录”\nref{S0409}；广东外语外贸大学2026年办法规定“\textbf{有违法违纪受处分记录者不予推荐}”\nref{S0134}。
\end{itemize}

\textbf{二、拿到资格之后再受处分，会被取消}

中南大学材料学院2027年细则明确“学生获得推免生资格后出现下列情况之一者，取消其推免生资格：……（二）\textbf{受刑事、行政、纪律处分的学生}”\nref{S0273}。南京工业大学2024年通知也规定“已经获得推免资格，但在毕业前因违反校纪、校规而受到学校相关处理的学生，学校将取消其推免生资格”\nref{S0527}。

\textbf{三、上位依据}。教育部令第41号第二十条规定，学校应记录学生学业、学术、品行方面的诚信信息，建立对失信行为的约束和惩戒机制，“对有严重失信行为的，可以规定给予相应的纪律处分，对违背学术诚信的，可以对其获得学位及学术称号、荣誉等作出限制”\nref{S0040}；第十九条规定无故缺席、违反考核纪律的可给予相应纪律处分\nref{S0040}。

\textbf{结论}：先去本院细则里找“处分期内”还是“无任何记录”这个措辞。前者你有机会，后者需要确认处分是否已被撤销或消除。
\nref{S0533}
\nref{S0248}
\nref{S0409}
\nref{S0134}
\nref{S0273}
\nref{S0527}
\nref{S0040}
\end{qa}

\iin{zsjyc（高校官网）}
\ith{违纪}
\ist{资格与排名}
\begin{qa}{M0103}{五级制成绩优良中怎么折算成保研用的分数}
\textbf{没有全国统一折算表，由各校自定；有明文折算表的学校多采用“优秀95、良好85、中等75、及格65、不及格55”这一类区间中值。}

可确认的完整实例只有一个（公开文件里极少见）：成都大学商学院依据《成都大学课程成绩管理办法》（成大教〔2024〕62号），在推免学业成绩计算中规定，若某课程原始成绩为五级制，按下列规则转换为百分制后再纳入加权计算\nref{S0079}：

\par\vspace{2pt}\nobreak\begin{center}\footnotesize
\begin{tabularx}{\textwidth}{>{\raggedright\arraybackslash}X>{\raggedright\arraybackslash}X>{\raggedright\arraybackslash}X>{\raggedright\arraybackslash}X>{\raggedright\arraybackslash}X>{\raggedright\arraybackslash}X}
\toprule
\textbf{五级制成绩} & \textbf{优秀} & \textbf{良好} & \textbf{中等} & \textbf{及格} & \textbf{不及格} \\ \midrule
百分制成绩 & 95 & 85 & 75 & 65 & 55 \\
\bottomrule
\end{tabularx}
\end{center}\vspace{2pt}\par

\textbf{为什么各校不同}：教育部令第41号第十三条只规定“考核分为考试和考查两种。考核和成绩评定方式，以及考核不合格的课程是否重修或者补考，\textbf{由学校规定}”\nref{S0040}，把成绩评定方式交给学校。

\textbf{常见替代处理}（不折算，直接不纳入）：
\begin{itemize}
  \item 温州医科大学规定“劳动教育、军事理论与实训、三段式考试课程实行二级制，\textbf{不纳入GPA统计}”\nref{S0563}；
  \item 大连理工大学商学院规定“课程属性为必修的导论课不纳入成绩计算范围”\nref{S0549}。
\end{itemize}

\textbf{结论}：五级制课程要么按本校折算表换成百分制计入，要么整门不纳入。两个方向对排名的影响可能相反（折算成95分会拉高，不纳入则不影响）。请向教务处索取本校的《课程成绩管理办法》原文。
\nref{S0079}
\nref{S0040}
\nref{S0563}
\nref{S0549}

申请时不要直接把其他学校的“优秀=95”等表格套用到自己：先查本校成绩管理办法是否有五级制折算，再确认目标院校要求的是课程成绩、绩点还是综合排名。提交时保留原始成绩单，并在附件说明折算来源和算法；如果校规没有公开口径，向教务部门取得书面说明，避免自行换算后与盖章排名不一致。
\end{qa}

\iin{成都大学商学院}
\ith{成绩折算}
\ist{资格与排名}
\begin{qa}{M0104}{P/NP的课算不算入排名}
\textbf{绝大多数学校不把P/NP（二级制、通过/不通过）课程计入排名——因为没有可加权的百分制分数。但这是学校自主事项，不是全国规定。}

\textbf{制度依据}：教育部令第41号第十三条规定“学生应当参加学校教育教学计划规定的课程和各种教育教学环节的考核，考核成绩记入成绩册”，同时“考核分为考试和考查两种。考核和成绩评定方式，以及考核不合格的课程是否重修或者补考，\textbf{由学校规定}”\nref{S0040}。P/NP属于“考查/二级制”评定方式，如何参与排名由学校定。

\textbf{各校实际写法}：

\begin{itemize}
  \item \textbf{明确不纳入}。温州医科大学2027年通知规定：“劳动教育、军事理论与实训、三段式考试课程实行\textbf{二级制，不纳入GPA统计}”\nref{S0563}。
  \item \textbf{按课程属性排除}。大连理工大学商学院规定“课程属性为必修的导论课不纳入成绩计算范围；截至公布综合成绩排名前一日，未出一考成绩的课程不纳入成绩计算范围”\nref{S0549}。
  \item \textbf{按学分门槛约束}。成都大学商学院要求“修读并获得本专业人才培养方案规定的前三学年（五年制为前四学年）\textbf{全部应修学分}”\nref{S0079}——P/NP课程虽然不计分，但必须“通过”才能满足学分条件。
\end{itemize}

\textbf{结论}：两件事要分开看——\textbf{计分}（P/NP一般不计入加权）和\textbf{学分}（P/NP课程通常计入“应修学分”，不及格/未通过会影响毕业与资格）。所以P/NP课不能放弃，只是考好考坏对排名无影响。

\textbf{注意}：有的学校写成“二级制课程不纳入GPA统计”，有的写“不纳入成绩计算范围”，措辞不同但效果接近。
\nref{S0040}
\nref{S0563}
\nref{S0549}
\nref{S0079}
\end{qa}

\iin{中华人民共和国教育部}
\ith{成绩折算}
\ist{资格与排名}
\begin{qa}{M0106}{第二学位在读能不能保研}
\textbf{不能。第二学士学位在读不在推免对象范围内，各校实施办法几乎一律排除。}

\textbf{学校层面的统一表述}：

\begin{itemize}
  \item 北京大学《推荐优秀应届本科毕业生免试攻读研究生工作实施办法》申请条件第（1）项：“纳入国家普通本科招生计划录取的应届毕业生（\textbf{不含专升本、第二学士学位学生}）”\nref{S0554}\nref{S0553}；
  \item 广东外语外贸大学2026年实施办法第二条：“本办法适用于学校全日制应届本科毕业生。应届，是指每年研究生报名开始时已进入毕业学年学习的本科生（\textbf{不含辅修学位、专升本、第二学士学位学生、本硕项目的国际班学生和独立学院学生}）”\nref{S0134}；
  \item 中南大学材料科学与工程学院2027年细则必备条件第1条：“纳入国家普通本科招生计划录取的应届生（\textbf{不含第二学士学位}）”\nref{S0273}；
  \item 成都大学商学院细则：“推荐对象为本校通过普通高考（含保送方式）招生录取或面向港澳台华侨学生招生录取的全日制应届本科毕业生（\textbf{不含专升本、第二学士学位学生}及国家政策不予推免的学生）”\nref{S0079}。
\end{itemize}

\textbf{上位依据}：教育部教学〔2006〕14号第十四条把推荐对象限定为“纳入国家普通本科招生计划录取的应届毕业生”\nref{S0454}；“免试”的定义是应届本科毕业生不必经过全国统考初试直接进入复试，指向的是本科应届毕业生。

\textbf{结论}：第二学士学位学生不具备申请推免的资格，只能走全国统考。另外，报考法律（法学）等专业学位时，第二学士学位另有专门规定——《2027年全国硕士研究生招生工作管理规定》第十六条明确“报考法律（法学）专业学位的，须符合……且报名前所学专业为法学专业（\textbf{获得法学第二学士学位的人员可报考}）”\nref{S0442}。

\textbf{注意}：不要与“辅修”混淆。辅修是主修之外的附加修读，主修学位仍是普通本科；第二学士学位是独立的学历学位层次。
\nref{S0554}
\nref{S0553}
\nref{S0134}
\nref{S0273}
\nref{S0079}
\nref{S0454}
\nref{S0442}
\end{qa}

\iin{北京大学}
\ith{第二学位}
\ist{资格与排名}
\section{交换学分}
\begin{qa}{M0065}{选修课算不算进保研排名？专业选修和通识选修要分开算吗}
\textbf{核心结论：选修课算不算、怎么算，各校差异极大，必须看本院细则的“课程计算清单”。} 五种真实写法并列：

\par\vspace{2pt}\nobreak\begin{center}\footnotesize
\begin{tabularx}{\textwidth}{>{\raggedright\arraybackslash}X>{\raggedright\arraybackslash}X}
\toprule
\textbf{学校（依据年份）} & \textbf{选修课口径} \\ \midrule
中南大学材料学院（2027细则） & 学业成绩＝“在校学习期间\textbf{所有教学环节（含全校性选修课）}加权平均裸分”，即选修课计入\nref{S0273} \\
南开大学（2021办法） & 排名依据是“\textbf{必修课（通识必修课、大类基础课、专业必修课）}平均学分绩”，只算必修类；且“计算学分绩所包含的课程范围……由各学院自主设定”\nref{S0561} \\
成都大学商学院（2027届） & 学业成绩＝“本专业培养方案中\textbf{所有必修课程}成绩的加权平均分”，选修课不计入\nref{S0079} \\
北京航空航天大学计算机学院（2027方案） & 加权平均分＝Σ(课程成绩×课程学分)/Σ课程学分，“列入计算范围的课程详见\textbf{课程计算清单}”\nref{S0579} \\
大连理工大学商学院 & 学习成绩只算必修课；但把\textbf{专业选修课}单列为另一项指标“专业知识”，算法为 专业知识成绩＝G×20/100，其中 G 为前三学年专业选修课程加权平均分\nref{S0549} \\
\bottomrule
\end{tabularx}
\end{center}\vspace{2pt}\par

可以看到三种完全不同的处理思路：\textbf{全算}、\textbf{只算必修}、\textbf{必修计分＋选修单列成另一项}。

\textbf{“通识选修”和“专业选修”要不要分开算}：多数学校不分开，而是按“是否属于培养方案内课程”这一条线切。大连理工的做法是例外——它把专业选修单独作为20分的一项，与“综合素质”“研究潜力”“研究能力”并列\nref{S0549}。这等于给了选修课一个固定的分值上限，考得再好也封顶。

\textbf{关键}：南开、北航都写明课程范围由学院划定并附清单，说明这件事的最终解释权在学院，不在学校。\textbf{去向学院教学办索取本院当年的“推免课程计算清单”，那才是唯一答案。}
\nref{S0273}
\nref{S0561}
\nref{S0079}
\nref{S0579}
\nref{S0549}
\end{qa}

\iin{中南大学}
\ith{课程范围}
\ist{资格与排名}
\begin{qa}{M0066}{体育课成绩算保研排名吗？我体育只有70分会不会拖后腿}
\textbf{看本校：多数学校体育课计入，但也有学校只把“体育课合格”设为门槛而不计入排名。}

\textbf{体育课“只当门槛”的写法}：
\begin{itemize}
  \item 北京科技大学2027年推免通知的基本条件第6条是“身心健康，\textbf{体育课成绩合格}”\nref{S0233}——即只要合格，具体分数不影响排名。
  \item 中山大学2026年强基计划转段方案要求“国家规定的\textbf{学生体质健康测试达标（免测除外）}”\nref{S0248}。
  \item 安徽师范大学2027年方案把“体质健康测试成绩”列入复核事项和公示内容，由体委负责核定\nref{S0646}。
\end{itemize}

\textbf{体育课“计入排名”的写法}：
\begin{itemize}
  \item 中南大学材料学院：学业成绩包含“在校学习期间\textbf{所有教学环节（含全校性选修课）}加权平均裸分”\nref{S0273}——体育课属教学环节，计入。
  \item 广东外语外贸大学2026年办法：要求“获得本专业本科人才培养方案应修总学分的80\%，\textbf{平均学分绩点不低于3.0}，无课程不及格记录”\nref{S0134}——体育课计入GPA。
  \item 成都大学商学院：学业成绩＝“本专业培养方案中\textbf{所有必修课程}成绩的加权平均分”\nref{S0079}——体育课属必修课，计入。
  \item 大连理工大学商学院：学习成绩只算必修课程\nref{S0549}。
\end{itemize}

\textbf{70分的影响有多大}：取决于两个因素——体育课在你校的实际学分（通常1—4学分），以及“计入还是只设门槛”。
\begin{itemize}
  \item 若只设门槛：70分只要及格，\textbf{完全不影响排名}。
  \item 若计入：设体育课共4学分、纳入计算的课程总学分180学分，那么体育70分相对80分的差距对总加权分的下拉约为 (70-80)×4/180 ≈ \textbf{-0.22分}——在名次密集处足以掉1—3名（参考湖南大学计算机学院公示的相邻名次分差常在0.1—0.2分\nref{S0131}）。
\end{itemize}

\textbf{结论}：先确认你校属于哪一类。如果只设门槛，就不用为体育分焦虑；如果计入，各校文件里没有“体育拖后腿”的豁免条款。
\nref{S0233}
\nref{S0248}
\nref{S0646}
\nref{S0273}
\nref{S0134}
\nref{S0079}
\nref{S0549}
\nref{S0131}
\end{qa}

\iin{北京科技大学}
\ith{课程范围}
\ist{资格与排名}
\begin{qa}{M0084}{转专业的学生保研排名怎么算？原来专业的课还算吗}
\textbf{主流做法：按现专业课程范围计算，原专业不在范围内的课不计入；名额也随现专业。}

\textbf{各校明文写法}：

\begin{itemize}
  \item \textbf{大连理工大学商学院}（最明确）：“\textbf{转专业学生按目前所学专业对应的必修课程范围计算，若所修课程不在规定课程范围内，则该课程不参与计算}”\nref{S0549}。也就是说，转专业后你的计算范围被重置为现专业的必修课清单，原专业修的、不在现专业培养方案里的课一律不进加权。
  \item \textbf{南开大学}（对交流/转学的写法）：“转学和校外交流学生可将学校\textbf{已经认定的学分}计入总学分”，但必修课平均学分绩“只计算学生在我校修读课程的成绩”\nref{S0561}。即“学分”可以认，“绩点”只算本校修读的。
  \item \textbf{成都大学商学院}：学业成绩＝“本专业培养方案中所有必修课程成绩的加权平均分”\nref{S0079}——按现专业培养方案算。
  \item \textbf{南京工业大学}：推荐对象明确“含\textbf{校内转专业学生}”\nref{S0527}，即转专业学生照常享有申请推免的资格，与“算哪几门课”是两个层面的事。
\end{itemize}

\textbf{结论}：分三个问题问清楚——
\begin{enumerate}
  \item \textbf{资格}：我转专业后还算是推免对象吗？（多数学校：是，如南京工业大学\nref{S0527}）
  \item \textbf{课程}：我按哪个专业的培养方案算加权？（多数学校：现专业，如大连理工\nref{S0549}、成都大学\nref{S0079}）
  \item \textbf{补修课}：转专业后补修的课程算不算？（各校不同；大连理工的写法是“在规定课程范围内就算”\nref{S0549}）
\end{enumerate}

\textbf{注意一处例外}：内蒙古大学把推免对象限定为“纳入国家普通本科招生计划录取的应届毕业生（\textbf{不含辅修学士学位学生、独立学院学生}）”\nref{S0544}——如果你同时是辅修学士学位身份，需要单独确认这一点。
\nref{S0549}
\nref{S0561}
\nref{S0079}
\nref{S0527}
\nref{S0544}
\end{qa}

\iin{大连理工大学}
\ith{转专业}
\ist{资格与排名}
\begin{qa}{M0085}{我大二去交换了一学期，交换的成绩算保研排名吗}
\textbf{两种相反做法，各校差异极大——有的学校完全不计入，有的学校认定后计入。}

\textbf{做法A：交换成绩一律不计入}（大连理工大学商学院，写得最直白）
“\textbf{交流生在交流期间原始成绩和转换后的成绩均不计入推免计算，含国内（包括港澳台地区）和国外高校}”\nref{S0549}。也就是“交换期间的成绩不进排名”，跟没修过一样。

\textbf{做法B：认定后可以计入}（南开大学）
“在推荐遴选时已经按照正常教学计划修满必修课程，\textbf{转学和校外交流学生可将学校已经认定的学分计入总学分}”，但必修课平均学分绩“\textbf{只计算学生在我校修读课程的成绩}”\nref{S0561}。也就是说，学分能认、能凑够毕业和资格所需的学分，但绩点只按本校修读课程算。

\textbf{做法C：按本校管理办法个案认定}
\begin{itemize}
  \item 西南交通大学（支教团要求）：“参加合作学校交流学习的，其免研资格课程由学院根据《西南交通大学本科生国（境）外交流学习管理办法》进行\textbf{合理认定}”\nref{S0061}。
  \item 内蒙古大学：“按照专业培养方案前三学年应修的必修课均已修完（\textbf{经学校批准赴境内外高校交流学习的学生按实施方案执行}）”\nref{S0544}。
\end{itemize}

\textbf{结论}：交换一学期对排名的可能影响是——完全不算（大连理工型）、或只认学分不算绩点（南开型）、或个案认定（西南交大、内大型）。

\textbf{给你的动作}：在出发交换之前，就书面问教务处三件事——① 交换课程能否认定为现专业培养方案内的课程；② 认定后是计入加权还是只计学分；③ 认定材料的提交截止时间。回来之后再问通常已经过了推免的成绩冻结时点。
\nref{S0549}
\nref{S0561}
\nref{S0061}
\nref{S0544}
\end{qa}

\iin{大连理工大学}
\ith{交换学分}
\ist{资格与排名}
\begin{qa}{M0087}{降级转专业后跟下一届一起排名，吃亏吗}
\textbf{这个问题没有公开规则可查——现行文件中未见“降级转专业后如何排名”的专门规定，属院校特定。} 但可以给你一个判断框架，避免凭感觉下结论。

\textbf{第一步：先确认你的学籍年级}。教育部令第41号第二十一条规定转专业由学校制定具体办法，“建立公平、公正的标准和程序，健全公示制度”\nref{S0040}；第二十五条规定学生应在最长学习年限内完成学业，休学次数和期限由学校规定。降级转专业意味着你的学籍年级下移，通常随之进入新年级的排名与名额体系。

\textbf{第二步：确认基数怎么算}。南开大学写明“以\textbf{排序时在校正式学籍学生数}为基准”\nref{S0561}，成都大学写“本专业\textbf{推荐对象基数}”\nref{S0079}——两者都指向“按排名时点的在籍人数”，而不是入学招生人数。你降级后进入的是下一届的池子。

\textbf{第三步：确认原专业成绩怎么处理}。大连理工大学商学院的口径是“转专业学生按\textbf{目前所学专业对应的必修课程范围}计算，若所修课程不在规定课程范围内，则该课程不参与计算”\nref{S0549}。即原专业的课多半不进加权。

\textbf{“吃不吃亏”取决于三个变量}：
\begin{enumerate}
  \item \textbf{两个年级/专业的竞争强度差}。降级后你从“原专业原年级”换到“新专业低一年级”，两个池子的强度可能差很多。
  \item \textbf{你已修课程的认定比例}。认定得多，你的加权样本更完整；认定得少，计入的课少，方差更大。
  \item \textbf{当年名额数}。名额按学院切，下一届名额若减少，即使你相对位置不差也会受影响。
\end{enumerate}

\textbf{结论}：没有“降级一定吃亏”或“一定占便宜”的通则。\textbf{最有效的做法是向教务处和学院书面索取三份数据}：你按哪一届/哪个专业排名、分母是多少、原专业课程如何认定。拿到这三份书面口径，你才能自己算出实际位置。
\nref{S0040}
\nref{S0561}
\nref{S0079}
\nref{S0549}
\end{qa}

\iin{中华人民共和国教育部}
\ith{转专业}
\ist{资格与排名}
\begin{qa}{M0105}{辅修双学位的成绩算不算保研排名}
\textbf{辅修成绩一般不进主修排名；排名按主修专业算。但内蒙古大学有一处需要特别注意的排除条款。}

\textbf{各校/上位文件的写法}：

\begin{itemize}
  \item \textbf{教育部令第41号第十六条}：“学生根据学校有关规定，可以申请辅修校内其他专业或者选修其他专业课程……学生修读的课程成绩（学分），\textbf{学校审核同意后，予以承认}。”\nref{S0040} 第三十六条：“对完成本专业学业同时辅修其他专业并达到该专业辅修要求的学生，由学校发给\textbf{辅修专业证书}。”\nref{S0040}\nref{S0366} ——辅修是主修之外的独立证书体系。
  \item \textbf{南京大学商学院}（接收端）：“当学生有\textbf{主修和辅修专业}时，\textbf{以主修专业为认定依据}”\nref{S0282}；在学科条件判断时也写作“本科专业（\textbf{不含辅修专业}）”\nref{S0282}。这从接收侧印证了“主修为准”。
  \item \textbf{内蒙古大学}（需要特别注意）：推免对象为“纳入国家普通本科招生计划录取的应届毕业生（\textbf{不含辅修学士学位学生、独立学院学生}）”\nref{S0544}——如果修读的是正式的“辅修学士学位”，按这条写法可能被排除在推免对象之外。这是少数学校的严格写法，务必核对本校条款。
  \item \textbf{广东外语外贸大学}：排除的是“辅修学位、专升本、第二学士学位学生、本硕项目的国际班学生和独立学院学生”\nref{S0134}。
\end{itemize}

\textbf{结论}：常规情况下，辅修成绩\textbf{不纳入主修排名}（辅修证书独立）。但有两类学校会把它变成问题：
\begin{enumerate}
  \item 把“辅修学士学位学生”明确排除在推免对象之外的（如内蒙古大学\nref{S0544}、广东外语外贸大学\nref{S0134}）；
  \item 把辅修课程认定为“本专业培养方案内课程”的。
\end{enumerate}

\textbf{给你的动作}：先确认你修的是“辅修专业证书”还是“辅修学士学位”——这两者在本校条款里的待遇可能完全不同。
\nref{S0040}
\nref{S0366}
\nref{S0282}
\nref{S0544}
\nref{S0134}
\end{qa}

\iin{中华人民共和国教育部}
\ith{辅修双学位}
\ist{资格与排名}
\begin{qa}{M0107}{辅修挂科会不会影响主修保研}
\textbf{一般不直接影响；但有两种学校会让它变成问题。}

\textbf{先说为什么不影响}。各校推免条件里说的“无不及格记录”，通常限定在\textbf{本专业培养方案内}的课程：
\begin{itemize}
  \item 成都大学商学院：成绩条件是“本专业培养方案中\textbf{所有必修课程}成绩的加权平均分……”，学分条件是修满“本专业人才培养方案规定的前三学年（五年制为前四学年）\textbf{全部应修学分}”\nref{S0079}；
  \item 广东外语外贸大学2026年办法：“获得本专业本科人才培养方案应修总学分的80\%，平均学分绩点不低于3.0，\textbf{无课程不及格记录}”\nref{S0134}；
  \item 内蒙古大学：“按照专业培养方案前三学年应修的必修课均已修完……\textbf{所修培养方案内课程均已合格}”\nref{S0544}。
\end{itemize}

辅修课程不在本专业培养方案内，所以按这些条款不算。

\textbf{两种例外}：

\begin{enumerate}
  \item \textbf{学校把“辅修学士学位学生”排除在推免对象外}。内蒙古大学实施办法第九条第（一）项：推免对象为“纳入国家普通本科招生计划录取的应届毕业生（\textbf{不含辅修学士学位学生、独立学院学生}）”\nref{S0544}；广东外语外贸大学2026年办法第二条同样排除“辅修学位”学生\nref{S0134}。这种情况下，问题不在辅修挂科，而在你的身份本身。
  \item \textbf{辅修课程被认定为“本专业培养方案内课程”}。这取决于本校的课程认定规则，需要个案确认。
\end{enumerate}

\textbf{上位依据}：教育部令第41号第三十六条规定辅修达到要求者由学校发给“辅修专业证书”\nref{S0040}\nref{S0366}——辅修是与主修学位并行的独立体系。

\textbf{结论}：辅修挂科一般不影响主修推免。\textbf{但如果你修的是“辅修学士学位”，请先核对本校条款有没有把它列入排除范围。}
\nref{S0079}
\nref{S0134}
\nref{S0544}
\nref{S0040}
\nref{S0366}
\end{qa}

\iin{成都大学商学院}
\ith{辅修双学位}
\ist{资格与排名}
\begin{qa}{M0108}{交换期间修的课如果没认定学分，排名怎么算}
\textbf{未认定的交换课程通常不计入排名；但更危险的是它会卡住“修满应修学分”这条资格条件。}

\textbf{一、计分：未认定＝不计入}
\begin{itemize}
  \item 大连理工大学商学院规定“交流生在交流期间原始成绩和转换后的成绩\textbf{均不计入}推免计算”，并且“\textbf{截至公布综合成绩排名前一日，未出一考成绩的课程不纳入成绩计算范围}”\nref{S0549}——两个条款叠加，未认定的交换课既没分也没成绩，直接不进加权。
  \item 南开大学的处理是“转学和校外交流学生可将学校\textbf{已经认定的学分}计入总学分”，但平均学分绩“只计算学生在我校修读课程的成绩”\nref{S0561}——即认定解决的是“学分”，不解决“绩点”。
\end{itemize}

\textbf{二、学分：未认定可能直接卡资格}
\begin{itemize}
  \item 成都大学商学院要求“修读并获得本专业人才培养方案规定的前三学年（五年制为前四学年）\textbf{全部应修学分}”\nref{S0079}；
  \item 云南师范大学2027年通知要求“学生须通过专业人才培养方案规定的\textbf{应修全部课程}”\nref{S0568}；
  \item 内蒙古大学要求“按照专业培养方案前三学年应修的必修课\textbf{均已修完}……所修培养方案内课程均已合格”\nref{S0544}。
\end{itemize}
只要本校是这种“修满型”表述，未认定的交换课程就意味着你的应修学分有缺口，可能直接被卡在资格审查上。

\textbf{三、认定有硬截止时点}
中南民族大学2027年通知写得最具体：“学业成绩的有效时间截止到2026年9月4日前（\textbf{从9月4日至推免工作结束这段时间，不处理本科生成绩录入、置换、认定等相关事宜}），各项奖励加分项目的有效时间截止到2026年9月4日前”\nref{S0594}。

\textbf{结论}：交换成绩未认定，对排名的影响是“不进加权”（通常中性或有利，因为不拉低平均）；但对资格的影响可能是致命的。\textbf{回来后第一件事就是办认定，不要拖到推免启动。}
\nref{S0549}
\nref{S0561}
\nref{S0079}
\nref{S0568}
\nref{S0544}
\nref{S0594}
\nref{S0061}\par\textbf{本版实操补充：} 先分别核对“课程是否计入培养方案”“成绩是否进入综合排名”两件事，二者可能不是同一口径。向本校教务部门提交交换成绩单、课程大纲、学分和学时证明，申请书面认定；在认定结果出来前，不要自行把外校成绩换算进排名。还要检查本校推免细则的成绩截止日、重修课口径及申诉时间，留下提交回执，避免只靠口头答复。

向教务部门问清并分别取得书面答复：课程是否计入毕业培养方案、是否纳入推免排名、采用原始成绩还是转换成绩、是否需要补修。若排名公示前仍未完成认定，要求说明暂不计入的依据及后续更正渠道；不要把课程“已修完”当作“已认定学分”。保留成绩单、课程大纲、学分学时和申请回执。
\end{qa}

\iin{大连理工大学}
\ith{交换学分}
\ist{资格与排名}
\begin{qa}{M0110}{保研名额出现并列排名怎么处理？有没有并列规则}
\textbf{不少学校有明文“同分排序规则”，次序通常是：先比加权成绩，再比科研成果/综合测评。}

\textbf{各校明文写法}：

\begin{itemize}
  \item \textbf{燕山大学}2027年通知：“符合遴选条件的学生按综合成绩……的顺序依次申请，且\textbf{综合成绩并列时，学习成绩高者优先}，额满为止。”\nref{S0533}
  \item \textbf{云南师范大学}2027年通知：“如出现综合成绩相同的情况，\textbf{依次按照专业平均成绩、科研成果成绩和实习实践成绩进行排名}。”\nref{S0568}
  \item \textbf{安徽师范大学}2027年方案：“普通类和硕师计划，如学生综合成绩相同，\textbf{按推免课程加权平均成绩高低排序}，如推免课程加权平均成绩仍相同，\textbf{按综合测评成绩高低排序}。”\nref{S0646}
\end{itemize}

\textbf{并列在真实名单里确实会发生}。湖南大学计算机学院2027届公示名单中，许家乐与邓广泽的学业成绩同为92.50、综合成绩同为92.50，并列标注为第4名（187人中）\nref{S0131}——说明该校的处理方式是“同位次并列”，而非强行分先后。

\textbf{两种处理风格}：
\begin{itemize}
  \item \textbf{并列同名次，后续名次跳号}（湖南大学\nref{S0131}）：同名次占用同一个序号，下一个名次从第6名开始。这种方式下，如果名额刚好卡在第4/第5之间，要靠学院另行裁决。
  \item \textbf{预先设定同分排序键}（燕山大学、云南师范大学、安徽师范大学\nref{S0533}\nref{S0568}\nref{S0646}）：在细则里先把比较顺序写死（加权成绩→科研成果→综合测评），出现同分直接按顺序判，不需要临时决策。
\end{itemize}

\textbf{结论}：如果名额恰好落在并列位置，先看本院细则有没有同分排序键。有，就按顺序判；没有，就应要求学院给出书面裁决依据——因为“临时决定谁上”在程序上最容易被申诉。

\textbf{注意}：同分排序键通常写在学院细则而不是校级文件里，所以要查本院文件，不能只看学校通知。
\nref{S0533}
\nref{S0568}
\nref{S0646}
\nref{S0131}
\end{qa}

\iin{zsjyc（高校官网）}
\ith{并列处理}
\ist{资格与排名}
\section{排名口径}
\begin{qa}{M0063}{保研排名到底按加权平均分还是学分绩点算？教务处和学院说法不一样}
\textbf{你的困惑是真实存在的：同一所学校里，校级“资格线”和院级“排序线”可能用两套口径。以学院《推免细则》为准。}

\textbf{三种主流口径}（都不是全国统一规定，教育部不作统一要求）：

\par\vspace{2pt}\nobreak\begin{center}\footnotesize
\begin{tabularx}{\textwidth}{>{\raggedright\arraybackslash}X>{\raggedright\arraybackslash}X}
\toprule
\textbf{学校（依据年份）} & \textbf{排序口径} \\ \midrule
中南大学材料学院（2027细则） & 学业成绩＝在校学习期间文化课（含全校性选修课）\textbf{加权平均裸分}，按专业排名；补考和重修合格按60分计\nref{S0273} \\
北京航空航天大学计算机学院（2027方案） & 综合成绩＝（\textbf{加权平均分}＋综合加分）/1.1；加权平均分＝Σ(课程成绩×课程学分)/Σ课程学分，原则上使用首次获得的课程成绩\nref{S0579} \\
成都大学商学院（2027届） & 学业成绩＝培养方案内所有必修课程的\textbf{加权平均分}；课程成绩采用首次获得的成绩\nref{S0079} \\
燕山大学（2027通知） & 综合成绩＝\textbf{加权平均学分绩点}＋创新创业竞赛获奖加分＋参军入伍加分\nref{S0533} \\
内蒙古大学 & 按学生所学课程\textbf{平均学分绩点}（保留两位小数）排名，不低于本专业前50\%\nref{S0544} \\
南开大学（2021办法） & 必修课（通识必修课、大类基础课、专业必修课）\textbf{平均学分绩}排名，前80\%以内；计算学分绩的课程范围与权重由各学院自主设定，权重不得低于50\%\nref{S0561} \\
\bottomrule
\end{tabularx}
\end{center}\vspace{2pt}\par

\textbf{为什么会出现“教务处和学院说法不一样”}：南开大学的写法最典型——校级只规定“平均学分绩”的计算底线（必修课、只算本校修读课程、前80\%），而“\textbf{计算学分绩所包含的课程范围及计算学分绩在综合成绩中所占权重由各学院根据本学院情况自主设定}”\nref{S0561}。教务处的口径是资格线，学院的口径是排序线，两者本就不同。

\textbf{怎么做}：向学院索取本院《推免细则》原文，找到“综合成绩”或“推免成绩”的计算公式和课程清单，以那份文件为准；口头答复不作数。
\nref{S0273}
\nref{S0579}
\nref{S0079}
\nref{S0533}
\nref{S0544}
\nref{S0561}
\end{qa}

\iin{中南大学}
\ith{排名口径}
\ist{资格与排名}
\begin{qa}{M0064}{算术平均分和加权平均分差0.3，会不会影响排名}
\textbf{会。而且在名额线附近，0.3分往往足以跨过好几个名次。}

先看一所学校的真实分差密度。湖南大学计算机学院2027届公示名单中，同一专业187人参与排名，相邻名次的分差是这样的\nref{S0131}：第1名94.42、第2名94.26（差0.16）、第3名93.44（差0.82）、第4名92.50、第4名92.50（并列）、第6名92.31（差0.19）、第11名91.40、第12名91.31（差0.09）、第13名91.11（差0.20）。\textbf{这意味着0.3分大约对应1—4个名次}；如果在名额线附近，就足以决定去留。

再看加权公式本身。北京航空航天大学计算机学院2027年方案规定：综合成绩＝（加权平均分＋综合加分）/1.1，加权平均分＝Σ(课程成绩×课程学分)/Σ课程学分\nref{S0579}。学分权重是直接的：一门4学分课与一门1学分课对总分的拉动相差4倍。所以“算术平均”和“加权平均”的差异主要来自课程学分分布——高分课程学分大时加权平均高于算术平均，反之相反。

\textbf{一个重要提醒}：本次检索到的各校官方文件中，\textbf{没有一所学校以“算术平均分”作为推免排序口径}。主流写法是“加权平均分”（中南大学\nref{S0273}、成都大学\nref{S0079}、北航\nref{S0579}）或“加权/平均学分绩点”（燕山大学\nref{S0533}、内蒙古大学\nref{S0544}）。如果你按算术平均自己算，很可能与学校口径不一致。\textbf{先确认本校公式，再算差值。}

顺便：中南大学的学业成绩之外还有专项评价加分，“各类型加分累计不超过10分（不含退役复学学生）”\nref{S0273}——0.3分在这里相对容易被加分弥补。
\nref{S0131}
\nref{S0579}
\nref{S0273}
\nref{S0079}
\nref{S0533}
\nref{S0544}
\end{qa}

\iin{湖南大学}
\ith{排名口径}
\ist{资格与排名}
\begin{qa}{M0068}{重修后的成绩是按重修分算还是按第一次算？会不会标注重修}
\textbf{记录上一定标注；计分上多数学校用“首次成绩”，少数按及格线（60分）计。}

\textbf{第一条：标注是法定义务。} 教育部令第41号第十八条明确：“学校应当健全学生学业成绩和学籍档案管理制度，真实、完整地记载、出具学生学业成绩，\textbf{对通过补考、重修获得的成绩，应当予以标注}。”\nref{S0040}\nref{S0366} 所以成绩单上的重修标记无法消除，接收院校能看到。

\textbf{第二条：计分方式各校不同，主流是两种写法。}

*写法A：用首次成绩*（即重修刷高不改变排名）
\begin{itemize}
  \item 成都大学商学院：“课程成绩计算\textbf{采用首次获得的成绩}”\nref{S0079}
  \item 北京航空航天大学计算机学院：“根据上级文件要求，\textbf{原则上使用学生首次获得的课程成绩}”\nref{S0579}
  \item 云南师范大学2027年通知：“学生学业成绩\textbf{原则上使用首次获得的课程成绩}”\nref{S0568}
  \item 燕山大学2027年通知：“加权平均学分绩点\textbf{按各科目首次考试成绩计算}”\nref{S0533}
  \item 大连理工大学商学院：“所有参评课程的成绩均\textbf{以一考成绩计算}”\nref{S0549}
\end{itemize}

*写法B：按合格线计分*
\begin{itemize}
  \item 中南大学材料学院2027年细则：“\textbf{补考和重修合格的课程成绩按60分计算}”\nref{S0273}
\end{itemize}

\textbf{结论}：无论哪一种，重修都不会“洗白”记录。写法A下重修对排名毫无帮助；写法B下重修后只能拿到60分，通常低于你原本可能考出的成绩，同样不划算。

\textbf{唯一值得重修的情形}：为了满足“无不及格记录”这类硬门槛（如广东外语外贸大学2026年办法要求“无课程不及格记录”\nref{S0134}）——补考通过即可满足，不必重修。
\nref{S0040}
\nref{S0366}
\nref{S0079}
\nref{S0579}
\nref{S0568}
\nref{S0533}
\nref{S0549}
\nref{S0273}
\nref{S0134}
\end{qa}

\iin{中华人民共和国教育部}
\ith{重修}
\ist{资格与排名}
\begin{qa}{M0073}{发一篇普刊能不能加分？会不会被认定为水刊反而扣分}
\textbf{普刊在多数学校根本进不了加分目录——不是“扣分”，而是“不加分”。“扣分”机制在本次检索的公开文件中未见。}

\textbf{全国层面的规则是“审核＋代表作评价”}：教育部推免规范管理文件要求，专家审核小组“对申请推免资格学生的科研创新成果、论文、竞赛获奖奖项及内容进行审核鉴定，\textbf{排除抄袭、造假、冒名及有名无实等情况}”，并“对学生提交的多篇科研成果实行\textbf{代表作评价}，评价重点聚焦到创新质量和个人贡献”；\textbf{未通过审核鉴定或答辩的，不得纳入推免遴选综合评价成绩计算体系}\nref{S0318}\nref{S0527}。

\textbf{各校对“什么论文能加分”的界定}：

\par\vspace{2pt}\nobreak\begin{center}\footnotesize
\begin{tabularx}{\textwidth}{>{\raggedright\arraybackslash}X>{\raggedright\arraybackslash}X}
\toprule
\textbf{学校（依据年份）} & \textbf{论文加分范围} \\ \midrule
内蒙古大学 & 仅限“在\textbf{核心期刊}上以独立作者或第一作者发表与学业相关的科研论文”，学分绩点加分最高不超过0.15\nref{S0544} \\
中南大学材料学院（2027细则） & 依托本校取得，第一作者：JCR 1区加0.6、2区加0.5、3区/4区/中文A类加0.4、EI和中文B类加0.3；多篇不累计取最高\nref{S0273} \\
安徽师范大学（2027方案） & 限该校《科研业绩分类和认定办法》“导向性期刊目录”\textbf{C2及以上}期刊\nref{S0646} \\
同济大学法学院（2027预报名） & “\textbf{收取版面费的期刊杂志上发表的论文不予认可}”\nref{S0242} \\
温州医科大学（2027通知） & 专家组“排除抄袭、造假、冒名及有名无实等情况”；“学生及其亲属或其他利益相关者合作的科研成果，原则上不纳入认定范围”\nref{S0562} \\
\bottomrule
\end{tabularx}
\end{center}\vspace{2pt}\par

\textbf{结论}：普刊通常连进入加分目录的资格都没有，所以不会有“扣分”。真正的风险是：如果你把普刊当作重要成果提交并申请加分，可能触发专家组的“有名无实”认定——后果是该项不计入加分，同时给审核组留下负面印象。\textbf{提交前先对照本校的期刊目录。}
\nref{S0318}
\nref{S0527}
\nref{S0544}
\nref{S0273}
\nref{S0646}
\nref{S0242}
\nref{S0562}
\end{qa}

\iin{shpg（高校官网）}
\ith{论文加分}
\ist{资格与排名}
\begin{qa}{M0086}{休学一年会不会影响保研资格和年级归属}
\textbf{休学本身不会取消推免资格，但会改变年级归属，通常意味着你要跟下一届一起排名、一起竞争。}

\textbf{学籍层面的规定}（教育部令第41号 第四节 休学与复学）\nref{S0040}\nref{S0366}：
\begin{itemize}
  \item 第二十五条：学生可以分阶段完成学业，除另有规定外，应当在最长学习年限（含休学和保留学籍）内完成学业；\textbf{休学次数和期限由学校规定}；
  \item 第二十七条：新生和在校学生应征参加中国人民解放军（含中国人民武装警察部队）的，\textbf{学校应当保留其入学资格或者学籍至退役后2年}；参加学校组织的跨校联合培养项目的，保留学籍；
  \item 第二十八条：休学学生应办理手续离校，休学期间学校为其保留学籍，但\textbf{不享受在校学习学生待遇}；
  \item 第二十九条：休学期满前应在学校规定的期限内提出复学申请，经学校复查合格方可复学。
\end{itemize}

\textbf{对推免的影响}：推免对象是“应届本科毕业生”（教学〔2006〕14号第十四条“纳入国家普通本科招生计划录取的应届毕业生”\nref{S0454}）。休学一年后复学，你的学籍年级随之下移，\textbf{当年不再属于应届毕业学年}，因此不参与当年的推免；复学后随新年级参与。这同时意味着名额池也变了——你进入的是下一届的名额分配方案（该校下一年的名额数、下一届的竞争强度）。

\textbf{一个例外是退役复学}。国家政策明确优待：教育部“退役大学生士兵计划”单列下达、专项使用\nref{S0052}；服役期间获三等战功、二等功以上奖励或二级以上表彰的可申请免初试攻读硕士\nref{S0442}；多校也把“参军入伍服兵役”纳入推免加分指标体系（如内蒙古大学\nref{S0544}）。

\textbf{结论}：休学不取消资格，但会让你“换一届重新排”。如果要休学，务必先向教务处书面确认复学后的年级归属和推免年次。
\nref{S0040}
\nref{S0366}
\nref{S0454}
\nref{S0052}
\nref{S0442}
\nref{S0544}
\nref{S0604}
\end{qa}

\iin{中华人民共和国教育部}
\ith{休复学}
\ist{资格与排名}
\begin{qa}{M0102}{保研排名里的学分绩点到底是4.0制还是5.0制，能直接比较吗}
\textbf{不能直接比较。绩点制各校不同（4.0制、5.0制、百分制混用），外校通常看的是“排名＋基数＋口径”，而不是绩点绝对值。}

\textbf{各校的绩点制实例}：

\begin{itemize}
  \item \textbf{4.0制}。湖南大学计算机学院2027届公示名单中，“学业成绩”与“GPA”两列并列：94.42分对应GPA 3.98，93.44对应3.97，92.50对应3.94，92.31对应3.95\nref{S0131}——典型的4.0制（且非简单线性换算）。
  \item \textbf{5.0制}。温州医科大学2027年通知把“必修课程及限选课程的平均学分绩点大于等于4.0”设为门槛\nref{S0563}\nref{S0562}——这个阈值只有在5.0制下才合理。
  \item \textbf{4.0制门槛}。广东外语外贸大学2026年办法要求“平均学分绩点不低于3.0（含3.0）”\nref{S0134}。
  \item \textbf{按绩点排名}。内蒙古大学规定“按学生所学课程平均学分绩点（保留到小数点后两位小数）排名，不低于本专业前50\%”\nref{S0544}；燕山大学用“加权平均学分绩点”\nref{S0533}。
  \item \textbf{直接用百分制}。中南大学、成都大学、北航等用加权平均分排序\nref{S0273}\nref{S0079}\nref{S0579}。
\end{itemize}

\textbf{外校怎么看}。北京大学力学与工程科学学院2026年夏令营的常见问题解答写得很清楚：“\textbf{不同学校的排名系统不同，按本校的排名规定填写}；如果本科院校学校只有'综合排名'或者只有'成绩排名'，那另外一项填写的时候，可以一样的排名，并在其它说明中注明……\textbf{总之，以院系盖章数据为准}”\nref{S0264}。四川大学网络空间安全学院、人工智能学院则要求“\textbf{排名需注明类型，如年级排名、班排名}”\nref{S0083}\nref{S0073}。

\textbf{结论}：填表时写清“百分比排名＋基数＋口径”三要素，比写绩点绝对值有用得多。
\nref{S0131}
\nref{S0563}
\nref{S0562}
\nref{S0134}
\nref{S0544}
\nref{S0533}
\nref{S0273}
\nref{S0079}
\nref{S0579}
\nref{S0264}
\nref{S0083}
\nref{S0073}
\end{qa}

\iin{湖南大学}
\ith{绩点换算}
\ist{资格与排名}
\begin{qa}{M0109}{我们学校的保研细则里前10\%是按什么基数的前10\%}
\textbf{基数通常写作“本专业（本年级）应届在籍学生数”，但三种操作会让分母变化，必须逐项确认。}

\textbf{各校对基数的明文写法}：

\begin{itemize}
  \item 成都大学商学院：“在本专业\textbf{推荐对象基数}的前20\%（含）”\nref{S0079}；
  \item 南开大学：“排名应在本专业本年级具有我校学籍学生的前80\%以内（\textbf{以排序时在校正式学籍学生数为基准}）”\nref{S0561}——注意“以排序时”这个措辞，意味着基数按排名时点的实际在籍人数算；
  \item 中国药科大学：国家级一流本科专业考生“本科推荐环节成绩排名前10\%”，其他本科专业前8\%\nref{S0418}——基数相同，按专业类别给不同比例；
  \item 中南大学材料学院：平行班和国际班学生“须在本专业应届生中排名前35\%”，卓越班学生“须在本专业\textbf{卓越班}应届生中排名前60\%”\nref{S0273}——同一学院内按班型分别算基数；
  \item 云南师范大学：“学院（学部）可按照\textbf{专业排名或班级排名}进行遴选，但\textbf{两种排名方式不得混用}”\nref{S0568}。
\end{itemize}

\textbf{会让分母变化的三个变量}：
\begin{enumerate}
  \item \textbf{是否含转专业、降级、休学复学学生}——按“排序时在校正式学籍学生数”算的（南开），这些人计入；按“应届毕业生数”算的可能不计入；
  \item \textbf{是否区分平行班与实验班}——如中南大学按班型分别算基数\nref{S0273}，等于实验班的竞争池被单独划出来；
  \item \textbf{用专业排名还是班级排名}——云南师范大学明确二者不得混用\nref{S0568}。
\end{enumerate}

\textbf{结论}：看到“前X\%”时，你要追问三个问题——分母是谁？按哪个时点算？用专业还是班级？\textbf{任何一个不同，前10\%的含义都不同。} 建议把这些写进书面问询，留底。
\nref{S0079}
\nref{S0561}
\nref{S0418}
\nref{S0273}
\nref{S0568}
\end{qa}

\iin{成都大学商学院}
\ith{排名口径}
\ist{资格与排名}
\begin{qa}{M0111}{候补的第11名有没有机会递补？什么时候能知道}
\textbf{有机会，但窗口很窄，而且“第11名”在多数学校进不了候补名单。关键是三个变量：本校候补上限、放弃人数、名额是否被学校收回。}

\textbf{候补名单通常有明确上限}：
\begin{itemize}
  \item 安徽师范大学2027年方案规定“普通类：学院分专业正式推荐人选依据推免综合成绩排名按1:1划定、\textbf{候补人选依据推免综合成绩排名划定2—3人}”\nref{S0646}——即候补只有2到3人，第11名进不去。
  \item 成都大学商学院写的是“候补人选人数按各专业推免名额数的\textbf{50\%}计算（计算结果不足1人的按1人计）”\nref{S0079}。
\end{itemize}

\textbf{递补的触发条件与规则}：
\begin{itemize}
  \item 南京林业大学风景园林学院2027年公示明确：“若已获推荐资格者因故放弃资格，或在公示期内被查实不符合推荐条件，其空余名额将依据综合成绩排序，\textbf{由后续候选人依次递补}”\nref{S0422}。
  \item 山东中医药大学中医学院2027年细则：“当排名靠前的学生放弃推免名额时，排名靠后的学生\textbf{依序递补，但不得超出相应成绩范围}。放弃推免名额的学生，须签订书面的放弃推免生名额承诺书”\nref{S0639}。
\end{itemize}

\textbf{放弃是真实且常见的}。湖南大学计算机学院2027届公示名单中，多名学生备注栏直接写着“放弃保研，留学”“放弃保研,就业”\nref{S0131}——这就意味着后面的同学有机会递补。

\textbf{什么时候能知道}：通常在校内公示期内到推荐名单报送前（9月上中旬）。如果学院名额没用完被学校收回重新分配（内蒙古大学、云南师范大学均如此规定\nref{S0544}\nref{S0568}），时间还会更晚，且名额可能不再回到原学院。

\textbf{结论}：第11名能否递补，取决于你所在专业/学院的候补上限和当年放弃的人数。\textbf{如果本校候补只划2—3人，你基本无望；如果候补按名额的50\%划且名额基数大，则有可能。} 最实际的动作是：公示期内紧盯名单变动，主动向学院说明“我在候补序列中，如有空缺请按序通知我”。
\nref{S0646}
\nref{S0079}
\nref{S0422}
\nref{S0639}
\nref{S0131}
\nref{S0544}
\nref{S0568}
\end{qa}

\iin{安徽师范大学}
\ith{递补}
\ist{资格与排名}
\begin{qa}{M0120}{我们学校把大一成绩也算进去，大二想翻盘还有机会吗}
\textbf{大一成绩基本无法剔除，但翻盘空间主要来自“加分项”，而不是刷分。}

\textbf{先确认一个事实：大一成绩几乎所有学校都算}。
\begin{itemize}
  \item 成都大学商学院：计算“前三学年（五年制为前四学年）”本专业培养方案内所有必修课程的加权平均分\nref{S0079}；
  \item 中南大学材料学院：须“完成专业培养方案前3年所有教学环节”\nref{S0273}；
  \item 大连理工大学商学院：要求“原则上修完前三学年培养方案要求的所有必修课程”\nref{S0549}；
  \item 南开大学：要求“已经按照正常教学计划修满必修课程”\nref{S0561}。
\end{itemize}

\textbf{关键变量是“学业成绩在综合成绩中的权重”}。权重越低，加分项能撬动的空间越大：
\begin{itemize}
  \item 南开大学校级规定：“计算学分绩权重\textbf{不得低于50\%}”\nref{S0561}——即加分项最多可占50\%；
  \item 温州医科大学：“\textbf{学业成绩在推免综合成绩中的权重不低于90\%}”\nref{S0563}——加分项最多只占10\%，翻盘空间很小；
  \item 成都大学商学院：“综合测评成绩＝学业成绩×90\%＋综合能力×10\%”（即学业占90\%）\nref{S0079}；
  \item 大连理工大学商学院：“T＝学习成绩×0.85＋思想素养(满分5分)＋综合素质和能力测评×0.1”\nref{S0549}；
  \item 中南大学材料学院：学业成绩满分100分＋专项评价成绩，专项加分“各类型累计不超过10分”\nref{S0273}。
\end{itemize}

\textbf{所以翻盘的两条路}：
\begin{enumerate}
  \item \textbf{靠加分项}——五类通用方向是参军入伍、志愿服务、国际组织实习、科研成果、竞赛获奖\nref{S0318}\nref{S0527}，各校自定目录与上限。如果本校学业权重只有50\%—85\%，加分项能真实改变名次。
  \item \textbf{靠后续课程}——大二大三的学分基数大、课程多，把后面每门课都考高，是唯一稳定的提分方式。
\end{enumerate}

\textbf{刷分是无效的路}：五所学校都按首次成绩计（成都大学、北航、云南师范大学、燕山大学、大连理工\nref{S0079}\nref{S0579}\nref{S0568}\nref{S0533}\nref{S0549}），重修不改善排名。

\textbf{结论}：大一成绩改不了，但先算清楚本校的“学业权重”再决定把精力放在刷绩点还是攒加分。
\nref{S0079}
\nref{S0273}
\nref{S0549}
\nref{S0561}
\nref{S0563}
\nref{S0345}
\nref{S0318}
\nref{S0527}
\nref{S0579}
\nref{S0568}
\nref{S0533}
\end{qa}

\iin{成都大学商学院}
\ith{排名口径}
\ist{资格与排名}
\section{公示异议}
\begin{qa}{M0078}{保研名额公示期一般几天？有异议该找谁}
\textbf{全国层面要求不少于7天；校内各环节通常还有3天左右的分段公示。}

\textbf{全国层面的硬规定}：教育部《全国普通高等学校推荐优秀应届本科毕业生免试攻读硕士学位研究生工作管理办法（试行）》（教学〔2006〕14号）第十五条第（五）项：“学校将审定的名单在各院（系）和校内网站公示，\textbf{公示期不少于7天}。对有异议的学生，学校要查明情况，公布处理结果……\textbf{未经公示的推免生资格无效}。”\nref{S0454} 教育部规范管理文件进一步要求“\textbf{公示期间公示内容不得随意变动，如有变动，须对变动部分做出说明，并对变动内容另行公示}”\nref{S0527}。

\textbf{各校实际公示期对照}：

\par\vspace{2pt}\nobreak\begin{center}\footnotesize
\begin{tabularx}{\textwidth}{>{\raggedright\arraybackslash}X>{\raggedright\arraybackslash}X}
\toprule
\textbf{学校（依据年份）} & \textbf{公示期} \\ \midrule
青海大学 & 院系公示3天，校领导小组审定后全校公示10天\nref{S0232} \\
南开大学（2021办法） & 全校公示\textbf{10个工作日}，未经公示的推荐资格无效\nref{S0561} \\
湖南大学 & 学院两轮排名公示各不少于3个工作日\nref{S0540}；计算机学院实际公示5天\nref{S0131} \\
燕山大学（2027） & 学院公示3天；全校名单公示7天\nref{S0533} \\
安徽师范大学（2027） & 学院不少于3天；学校不少于7天\nref{S0646} \\
天津大学（2027） & 拟录取名单公示7日\nref{S0520} \\
\bottomrule
\end{tabularx}
\end{center}\vspace{2pt}\par

\textbf{有异议该找谁}：① \textbf{学院推免生遴选工作小组/监督小组}——山东中医药大学中医学院要求“提交实名书面申诉材料至学院推免遴选工作小组”\nref{S0639}，青海大学、云南师范大学也要求学院组建监督小组受理异议\nref{S0232}\nref{S0568}；② 学校教务处/研究生院；③ 学校学生申诉处理委员会（教育部令第41号第五十九条\nref{S0040}）。

\textbf{注意}：“未经公示的推免生资格无效”是全国层面的硬约束，公示是资格成立的必要条件。
\nref{S0454}
\nref{S0527}
\nref{S0232}
\nref{S0561}
\nref{S0540}
\nref{S0131}
\nref{S0533}
\nref{S0646}
\nref{S0639}
\nref{S0520}
\nref{S0040}
\end{qa}

\iin{教育部学生服务与素质发展中心（学信网/研招网）}
\ith{公示异议}
\ist{资格与排名}
\begin{qa}{M0079}{公示排名跟自己算的不一样，能申诉吗？流程是什么}
\textbf{能申诉。三层渠道，越往上程序越正式。}

\textbf{第一层：学院复议（最快，必须在校名单报送前）}。湖南大学2027年工作安排：学生对推荐结果有异议的，\textbf{填写《湖南大学推免生复议申请表》（一式三份），向所在单位推荐工作小组申请复议}；工作小组审定并在表中\textbf{详细记录复议过程及结论，组长签字并加盖公章}；\textbf{“复议结论应在推荐名单报送推免生遴选工作领导小组之前告知申请人”}\nref{S0540}\nref{S0537}。这个“报送前告知”很关键——学院必须在名单上移前给你答复。

\textbf{第二层：书面实名申诉到学院推免工作小组/监督小组}。山东中医药大学中医学院2027年细则第十三条：“学院初选名单公示期间，学生如有异议，\textbf{提交实名书面申诉材料至学院推免遴选工作小组}”\nref{S0639}。云南师范大学要求“各学院（学部）应组建专门的推免工作监督小组（成员3或5人）……\textbf{受理推免工作中的异议和申诉}”\nref{S0568}。北京大学2027年推免办法要求“各院系应畅通学生咨询及申诉渠道……\textbf{及时调查处理并予以反馈}”\nref{S0553}。

\textbf{第三层：学校学生申诉处理委员会}（教育部令第41号第五十九至六十一条）\nref{S0040}：学生对学校的处理或处分决定有异议的，可在接到决定书之日起\textbf{10日内}提出书面申诉；申诉处理委员会应在接到书面申诉之日起\textbf{15日内}作出复查结论并告知申诉人，情况复杂可延长15日。

\textbf{操作建议}：① 一定书面，附自己算的口径、课程清单、依据条文；② 引条文——指出“未按本校细则第X条排序规则执行”或“公示内容变更未重新公示”（违反教育部“公示期间内容不得随意变动，如有变动须另行公示”\nref{S0527}）；③ 卡时间——复议要在名单报送前，申诉要在10日内；④ 用数据说话。
\nref{S0540}
\nref{S0537}
\nref{S0639}
\nref{S0568}
\nref{S0553}
\nref{S0040}
\nref{S0527}
\end{qa}

\iin{湖南大学}
\ith{公示异议}
\ist{资格与排名}
\begin{qa}{M0095}{退役复学后保研有什么优惠？士兵计划是加名额还是单独划线}
\textbf{“退役大学生士兵计划”是研究生招生环节的单列专项（走统考、单独划线），不是“给学院加一个推免名额”；另有“免初试”通道；退役复学在推免推荐阶段体现为加分或门槛放宽。}

\textbf{一、退役大学生士兵计划＝统考环节的单列专项。} 教育部办公厅2026年通知（教学厅〔2025〕8号）：“'退役大学生士兵计划'……\textbf{在全国硕士研究生招生计划总规模内单列下达、专项使用，不得调整为普通计划或其他专项计划}。招生单位自主确定该计划考生初试成绩要求（一般不高于本单位相应学科专业普通计划考生初试成绩要求）……同等条件下服役时间长的考生优先录取。”2026年由北京大学、清华大学等532所高校承担\nref{S0440}\nref{S0052}。参加该计划仍要参加初试。

\textbf{二、真正的“免初试”通道。} 《2027年全国硕士研究生招生工作管理规定》第二十五条：服役期间获得\textbf{三等战功、二等功以上奖励或者二级以上表彰}、符合报考条件的退役人员，可申请免初试攻读硕士研究生，须在教育部规定的统考报名时间内登录推免服务系统报名，招生单位组织复试，复试合格者予以拟录取\nref{S0442}\nref{S0051}。

\textbf{三、退役复学在推免推荐阶段的作用。} 教育部规范管理文件把“参军入伍服兵役”列为可纳入遴选指标体系的五类因素之一\nref{S0318}\nref{S0527}。中南大学材料学院2027年细则：退役复学学生学业成绩门槛为“本专业应届生中排名前60\%”（平行班为前35\%），且“各类型加分累计不超过10分（\textbf{不含退役复学学生}）”\nref{S0273}；内蒙古大学也把应征入伍退役学生纳入遴选指标体系\nref{S0544}。

\textbf{结论}：优惠落在三处——校内推免的加分/门槛放宽、退役大学生士兵计划（统考单列）、符合条件者免初试。路径不同，报名时间也不同。
\nref{S0440}
\nref{S0052}
\nref{S0435}
\nref{S0442}
\nref{S0051}
\nref{S0318}
\nref{S0527}
\nref{S0273}
\nref{S0544}
\end{qa}

\iin{教育部学生服务与素质发展中心（学信网/研招网）}
\ith{退役政策}
\ist{资格与排名}
\begin{qa}{M0114}{艺术团高水平运动队的保研通道跟普通学生一样吗}
\textbf{不一样。这类特长生的通道有明确限制——教育部规定不得单列名额、不得制定单独办法；有些学校只给“门槛放宽”，不给“加分”。}

\textbf{全国层面的约束}。教育部推免规范管理文件要求：“合理设置各遴选指标所占权重及单项指标上限分值，学生在某一方面中有多项加分情况时，原则上只取一项，\textbf{不得为仅符合单一或部分遴选指标的学生单列计划或破格推荐}”，并在括号里点名了“特殊学术专长者、\textbf{高水平运动员}、退役大学生士兵、竞赛获奖者等”\nref{S0527}。云南师范大学2027年通知也重申“\textbf{不得为仅符合单一或部分遴选条件的学生（如高水平运动员、退役大学生士兵、竞赛获奖者等）单列计划}”\nref{S0568}。

\textbf{更直接的一条}：湘潭大学2026年通知转述《教育部关于进一步加强普通高等学校高水平运动队建设管理的意见》（教体艺〔2022〕1号）：“\textbf{不得给高水平运动员单列推免名额或制定单独推免办法}，符合学校推免生申请条件的高水平运动员，可按照学校统一规定和办法申请推免资格。”\nref{S0567}

\textbf{实际存在的“倾斜”形式——放宽门槛，而非单列名额}。中南大学材料科学与工程学院2027年细则：平行班和国际班学生须“在本专业应届生中排名前35\%”，而“\textbf{退役复学学生、高水平运动员、高水平艺术团成员、参与社会工作志愿服务的懂'双语'少数民族学生须在本专业应届生中排名前60\%}”\nref{S0273}。即走同一套排名、同一个名额池，只是资格线更低。

\textbf{加分方面}。湖南大学2027年单独公示了“体育类、艺术类竞赛获国家级及以上奖励、志愿服务加分名单”\nref{S0541}；中南大学材料学院的加分细则也含“文体素质比赛获奖加分”（国家级一等奖排名第一加0.6分等）\nref{S0273}。

\textbf{结论}：艺术团、高水平运动队学生与普通学生\textbf{同池竞争}——差别在于资格门槛更低和文体类竞赛加分。\textbf{不存在“单独的保研名额”}，教育部明确禁止单列。
\nref{S0527}
\nref{S0568}
\nref{S0567}
\nref{S0273}
\nref{S0541}\par\textbf{本版实操补充：} 不要把“特长生身份”直接等同于单独推免资格。先查本校当年推荐对象范围、综合评价指标和是否有单列培养项目，再确认相关身份材料由哪个部门审核。全国规则与校级细则的适用层级不同；如果文件没有明确列出你的身份，不要按往年口头案例推断，向教务处或推免工作组提交书面问题并保存答复。

判断时把三件事分开查：是否属于当年可推荐对象、是否占普通推免名额、综合评价是否存在明确加分或单列培养项目。向负责审核的教务或体育艺术部门索要当年文件及适用对象定义；如果身份名称相近但文件未列明，先书面确认，不要以竞赛成绩或历史案例推定资格。
\end{qa}

\iin{教育部学生服务与素质发展中心（学信网/研招网）}
\ith{特长通道}
\ist{资格与排名}
\section{排名时点}
\begin{qa}{M0080}{保研是前六学期还是前五学期成绩？大三下成绩出得晚怎么办}
\textbf{校内推免用前六学期/前三学年；夏令营只能用前五学期。这就是“时点错配”的根源。}

\textbf{校内推免（前六学期/前三学年）}：
\begin{itemize}
  \item 成都大学商学院：要求修满“前三学年（五年制为前四学年）”全部应修学分，时点是“\textbf{推免工作当年春季学期正常期末考试成绩统一录入后}”\nref{S0079}；
  \item 中南大学材料学院：须“完成专业培养方案\textbf{前3年}所有教学环节”，“\textbf{成绩审核以截止9月10日下午17:00教务管理系统中公布的加权平均分和专业排名为准}”\nref{S0273}；
  \item 大连理工大学商学院：“截至\textbf{公布综合成绩排名前一日}，未出一考成绩的课程不纳入成绩计算范围”\nref{S0549}；
  \item 教育部教学厅〔2014〕5号：“推荐工作统一于每年的\textbf{9月25日前}结束”\nref{S0205}。
\end{itemize}

\textbf{夏令营/预推免（前五或前六学期）}：
\begin{itemize}
  \item 上海交通大学计算机学院2026年夏令营：“本科\textbf{前五或前六学期}成绩的专业排名在前20\%”，并需“前三学年（或前五学期）学习成绩专业排名证明”\nref{S0358}；船建学院2027年预推免要“\textbf{前五或前六个学期}”成绩单与排名证明\nref{S0575}；
  \item 合肥工业大学计算机与信息学院2026年夏令营：“本科\textbf{前五学期}加权成绩排名在本专业前列”\nref{S0098}；
  \item 北京师范大学社会学院2026年夏令营：“学生综合成绩（\textbf{本科前五学期成绩}）排名在本专业前30\%”\nref{S0326}；
  \item 兰州大学文学院2026年夏令营：“在校\textbf{前五学期}学业成绩排名在专业或班级前20\%”\nref{S0094}。
\end{itemize}

\textbf{怎么办}：多数夏令营在7月前后报名，第6学期成绩往往还没出齐，所以只能用前五学期排名——这是全国性现象，接收院校通常接受“前五学期”或“前六学期中已出部分”。

\textbf{结论}：7月夏令营用前五学期，9月校内推免用前六学期/前三学年。两者不是同一个名次，别拿夏令营的排名去推推免的结果。
\nref{S0079}
\nref{S0273}
\nref{S0568}
\nref{S0549}
\nref{S0205}
\nref{S0358}
\nref{S0575}
\nref{S0098}
\nref{S0326}
\nref{S0094}
\end{qa}

\iin{成都大学商学院}
\ith{排名时点}
\ist{资格与排名}
\begin{qa}{M0081}{我们学校大三下才出成绩，夏令营只能用前五学期排名吗}
\textbf{是的，多数夏令营只收前五学期排名，但普遍接受“暂缺后补”或“预估并备注”这两种补救方式。}

\textbf{接收院校的三种宽松写法}：

\begin{enumerate}
  \item \textbf{明确允许后补}。山东大学经济研究院2026年夏令营：申请材料包括“前五学期的成绩单原件与总评成绩排名证明（均需单位盖章），\textbf{如果成绩排名暂时没有，九月份补寄亦可}”\nref{S0026}。
  \item \textbf{接受“无三年成绩则用五学期”}。北京交通大学计算机科学与技术学院2026年夏令营要求“本科前三学年（\textbf{如无，则前五学期}）成绩单（教务部门盖章）”和“本科前三学年（如无，则前五学期）成绩排名证明”\nref{S0104}。
  \item \textbf{接受预估并备注}。中国科学院大学生命科学学院2026年夏令营写得很具体：“'年级（专业）排名'填写\textbf{前6学期排名情况（如第6学期成绩未出，可按前5学期排名填写）}；若无官方统一排名，可结合个人成绩情况\textbf{预估排名，并在材料中单独备注}'本专业无官方排名，此为个人预估'”\nref{S0077}。
\end{enumerate}

\textbf{“排名怎么填”这件事接收院校也给了指导意见}。北京大学力学与工程科学学院2026年夏令营常见问题解答：“不同学校的排名系统不同，\textbf{按本校的排名规定填写}；如果本科院校只有'综合排名'或者只有'成绩排名'，那另外一项填写的时候，可以一样的排名，并在其它说明中注明……\textbf{排名不确定如何填写？答：如没有准确名次，请结合自己的实际情况尽量准确预估}……总之，以院系盖章数据为准。”\nref{S0264}

\textbf{给你的三步操作}
\begin{enumerate}
  \item \textbf{先报，用前五学期口径}，并在申请表备注栏注明“截至第5学期”；
  \item \textbf{排名未出的，写“暂无官方排名，此为预估”}（中科院大学生命科学学院明确允许\nref{S0077}）；
  \item \textbf{9月拿到前六学期排名后向已入营/已联系的院校补寄}（山东大学经济研究院明确接受补寄\nref{S0026}）。
\end{enumerate}

另外注意，浙江大学医学院第一临床医学院2026年夏令营要求排名证明“\textbf{应包含所学专业同年级人数}”\nref{S0395}——所以材料里不仅要写名次，还要写基数。
\nref{S0026}
\nref{S0104}
\nref{S0077}
\nref{S0264}
\nref{S0395}
\nref{S0098}
\nref{S0376}
\end{qa}

\iin{山东大学}
\ith{排名时点}
\ist{资格与排名}
\begin{qa}{M0089}{强基计划的学生保研是走推免还是转段}
\textbf{强基生走“本研转段”，是独立通道，不与普通推免共用名额。}

\textbf{制度基础}：教育部《关于在部分高校开展基础学科招生改革试点工作的意见》（教学〔2020〕1号）建立强基计划\nref{S0045}\nref{S0368}；《教育部办公厅关于做好强基计划学生管理和研究生衔接培养有关工作的通知》（教学厅〔2020〕9号）安排衔接培养。各校据此制定转段方案，如中山大学参照《中山大学强基计划学生培养管理办法》制定本研转段实施方案\nref{S0248}。

\textbf{三个关键差异（都是“转段≠推免”的证据）}：

\begin{enumerate}
  \item \textbf{不占用推免指标，但占用导师招生计划}。中国科学技术大学2024年强基生转段通知：学位层次“须符合拟报考的研究生招生院系及导师2025年招生安排，\textbf{不占用其推免录取指标，但占用导师当年硕士或博士招生计划}（录取纳入教育部专项计划除外）”\nref{S0497}。
  \item \textbf{只能申请校内}。华中科技大学2026年转段接收通知：“获得转段资格的强基生，\textbf{只能申请校内相关院系的硕士或直接攻读博士}”\nref{S0536}。普通推免则可自主报考外校——这是最本质的区别。
  \item \textbf{独立的时间节点与考核}。湖南大学2027年工作安排里“强基转段学生工作（9月14日前）”与普通推免各项公示并列，由化学化工学院单独完成\nref{S0540}。中山大学的转段条件包括“国家规定的学生体质健康测试达标（免测除外）”“主修专业培养方案规定的公共必修课、专业必修课和专业选修课考核无不及格情况”“在校期间未因考试作弊或学术不端行为受到处分”\nref{S0248}。转段接收面试参照推免生接收程序，考核外语水平、专业知识、科研能力、学术潜力、思想政治素质和品德等\nref{S0536}。
\end{enumerate}

\textbf{直博生指标可能单列}：中山大学2026年转段方案明确“本年度强基生转段按照直博生录取，\textbf{博士研究生指标将单列下达}”\nref{S0248}。

\textbf{结论}：强基生走转段，不占推免名额、不能外校、单独考核；推免可外校、用推免名额。两者互不替代。
\nref{S0248}
\nref{S0536}
\nref{S0497}
\nref{S0045}
\nref{S0368}
\nref{S0540}
\end{qa}

\iin{中山大学}
\ith{强基计划}
\ist{资格与排名}
\begin{qa}{M0119}{保研排名的截止时间点是怎么定的？期末考完才截止吗}
\textbf{不是等期末考完统一裁，而是学校设一个“数据冻结时点”，晚于该时点的成绩可能不被采纳。}

\textbf{不同学校的冻结时点写法}：

\begin{itemize}
  \item \textbf{中南大学材料学院}（最明确）：“成绩审核以\textbf{截止9月10日下午17:00教务管理系统中公布的加权平均分和专业排名为准}。”\nref{S0273}
  \item \textbf{中南民族大学}：“学业成绩的有效时间\textbf{截止到2026年9月4日前}（从9月4日至推免工作结束这段时间，\textbf{不处理本科生成绩录入、置换、认定等相关事宜}），各项奖励加分项目的有效时间截止到2026年9月4日前。”\nref{S0594}
  \item \textbf{大连理工大学商学院}：“\textbf{截至公布综合成绩排名前一日，未出一考成绩的课程不纳入成绩计算范围}。”——即晚出分的课直接被排除，而不是等等看\nref{S0549}。
  \item \textbf{湖南大学}（时间表）：学院学业综合成绩排名公示7月中旬前、学院推免遴选总成绩排名公示8月底前\nref{S0540}\nref{S0537}。
  \item \textbf{安徽师范大学}：9月1日前完成已有课程成绩、综合测评、外语成绩、体质健康达标等情况的复核；9月2日至5日完成分专业推免课程成绩计算并排名\nref{S0646}。
\end{itemize}

\textbf{全国层面的框架}：教育部教学厅〔2014〕5号规定“推荐工作统一于每年的\textbf{9月25日前结束}”\nref{S0205}——所有学校的内部时点都必须早于这个上限。

\textbf{结论}：
\begin{itemize}
  \item \textbf{时点由学校自定}，通常在8月底至9月中，不是“期末考完自然截止”。
  \item \textbf{晚出分课程的处理有两种}：直接被排除（大连理工\nref{S0549}）、或该时段不办理录入（中南民族大学\nref{S0594}）。
  \item \textbf{“最后一次刷分”的机会窗口在时点之前}。如果你想靠某门课提分，必须在学校公布的成绩冻结日之前出分并录入。
\end{itemize}

\textbf{务必查}：本院细则里的“成绩审核截止时间”和“未出分课程处理方式”两条，这两条决定了你大三下学期能不能“补救”。
\nref{S0273}
\nref{S0594}
\nref{S0549}
\nref{S0540}
\nref{S0537}
\nref{S0646}
\nref{S0205}
\end{qa}

\iin{中南大学}
\ith{排名时点}
\ist{资格与排名}
\section{排名证明}
\begin{qa}{M0076}{保研排名证明谁开？教务处的章和学院的章哪个有效}
\textbf{两种章都可能有效，取决于目标院校当年的要求——不同学校写得不一样，必须逐校核对。}

\textbf{接受“学校或学院教务部门”的写法（较宽松，占多数）}：
\begin{itemize}
  \item 上海交通大学计算机学院2026年夏令营：成绩单与“本科前三学年（或前五学期）学习成绩专业排名证明”均“需加盖\textbf{所在学校或学院教务部门公章}”\nref{S0358}；凯原法学院写“加盖\textbf{所在学校教务处或所在学院教务办}的公章”\nref{S0240}。
  \item 上海财经大学经济学院2027年预报名：排名证明“由所在学校教务部门或学院填写\textbf{加盖公章}后扫描”\nref{S0138}。
  \item 四川大学网络空间安全学院/人工智能学院：成绩单和排名证明“需加盖\textbf{所在学校教务处或所在学院公章}”\nref{S0083}\nref{S0073}。
  \item 浙江大学医学院第一临床医学院：均“须加盖\textbf{学校教务处或院系公章}”\nref{S0395}。
  \item 东南大学计算机学院：“已有学校教务部门或院系盖章的成绩排名证明，申请表中学校教务部门或院系签章可免”\nref{S0126}。
\end{itemize}

\textbf{只认一种章的写法}：同济大学法学院写“本科成绩单(\textbf{由学校教务处盖章})”\nref{S0242}；浙江大学物理学院要求“年级排名证明（\textbf{需加盖院系公章}）”\nref{S0293}；兰州大学文学院要求“\textbf{需加盖学院公章}”\nref{S0094}。

\textbf{建议}：① 多数院校接受“学校或学院教务部门”二者之一，学院教务章在多数情况下有效；② 但存在只认学校教务处章的写法（同济大学法学院\nref{S0242}），也有只认院系章的（浙江大学物理学院\nref{S0293}）；③ \textbf{一次开两份（教务处版＋学院版）}，按目标院校要求提交，教务处的章覆盖面更广；④ 注意分辨“教务章”与“行政章”——要求写的是“教务部门公章”，学院行政公章可能不被接受。
\nref{S0358}
\nref{S0240}
\nref{S0242}
\nref{S0138}
\nref{S0083}
\nref{S0073}
\nref{S0395}
\nref{S0126}
\nref{S0293}
\nref{S0094}
\end{qa}

\iin{上海交通大学}
\ith{排名证明}
\ist{资格与排名}
\begin{qa}{M0077}{排名证明上写的是专业排名还是年级排名，外校怎么看}
\textbf{外校看的不是名次本身，而是“名次＋基数＋口径”。写专业排名还是年级排名都可以，但必须真实并注明。}

\textbf{接收院校的要求}：
\begin{itemize}
  \item 四川大学网络空间安全学院、人工智能学院：“\textbf{排名需注明类型，如年级排名、班排名}”\nref{S0083}\nref{S0073}；
  \item 浙江大学物理学院：材料清单写“\textbf{年级排名证明}（需加盖院系公章）”\nref{S0293}；
  \item 浙江大学医学院第一临床医学院：“总评成绩排名证明（\textbf{应包含所学专业同年级人数}…）”\nref{S0395}；
  \item 北京林业大学外语学院：对不同出身给不同门槛——“双一流”高校学生前五学期\textbf{专业排名前50\%}、外语类院校前30\%、其他学校校本专业\textbf{前15\%}\nref{S0023}。同一份“专业排名”，基数不同含义完全不同；
  \item 北京大学基础医学院：“因排名系统不同，请依照所在本科院校排名的规定填写…\textbf{请在申请表上备注说明}”\nref{S0301}。
\end{itemize}

\textbf{北京大学力学与工程科学学院FAQ给了最清楚的处理办法}：“不同学校的排名系统不同，\textbf{按本校的排名规定填写}；如果本科院校学校只有'综合排名'或者只有'成绩排名'，那另外一项填写的时候，\textbf{可以一样的排名，并在其它说明中注明}……总之，\textbf{以院系盖章数据为准}。”\nref{S0264}

\textbf{建议}：
\begin{enumerate}
  \item \textbf{名次必须配基数}——“第10名”没意义，“第10名/共187人”才有意义（湖南大学计算机学院公示就是这格式\nref{S0131}）。
  \item \textbf{注明口径}——年级/专业排名、含不含全部课程、截至第几学期。
  \item \textbf{只有一种排名时，另一种可填相同并备注}（北大FAQ明确允许\nref{S0264}）。
\end{enumerate}
\nref{S0083}
\nref{S0073}
\nref{S0293}
\nref{S0395}
\nref{S0023}
\nref{S0301}
\nref{S0264}
\nref{S0131}
\end{qa}

\iin{四川大学}
\ith{排名证明}
\ist{资格与排名}
\begin{qa}{M0116}{保研资格是有资格就够了还是要排名靠前才算}
\textbf{要分开三个概念：资格、名额、录取。你真正缺一不可的是前两个。}

\begin{enumerate}
  \item \textbf{推免资格}（“有没有资格”）：这是本科学校认定的，也是最稀缺的。教育部教学〔2006〕14号第十二条规定了名额比例——“教育部批准设立研究生院的高等学校一般按应届本科毕业生数的\textbf{15\%左右}确定；未设立研究生院的'211工程'建设高校一般按\textbf{5\%左右}；其他高等学校一般按\textbf{2\%}确定，其中初次开展推荐工作的前3年每年一般按\textbf{1\%}确定”\nref{S0454}。也就是说，“资格”不是人人有的，它本身就是稀缺资源。
\end{enumerate}

\begin{enumerate}
  \item \textbf{名额}（“排名够不够拿到指标”）：这才是“排名靠前才算”的部分。山东中医药大学中医学院2027年细则把申请条件写成“\textbf{综合测评成绩排名范围要求：前4学年综合测评总成绩须排在本专业年级的前30\%}”\nref{S0639}；云南师范大学2027年通知明确“\textbf{综合成绩低于60分的学生不予推荐}……按照综合成绩从高到低确定推荐资格名单”\nref{S0568}。所以：\textbf{满足基本条件 ≠ 拿到名额}，中间隔着排序。
\end{enumerate}

\begin{enumerate}
  \item \textbf{录取}（“目标院校要不要你”）：拿到名额后，还要通过接收院校的复试/考核。教育部教学厅〔2014〕5号规定“所有推免生均享有依据招生政策自主选择报考招生单位和专业的权利，所有推免名额（除有特殊政策要求的专项计划外），均可向其他招生单位推荐”\nref{S0205}——你有自主报考权，但能不能被接收是对方决定的。湖南大学计算机学院2027届公示名单里有人备注“放弃保研，留学”“放弃保研,就业”\nref{S0131}，也说明“拿到名额”和“最终入学”是两个环节。
\end{enumerate}

\textbf{结论}：
\begin{itemize}
  \item \textbf{排名靠前是拿到名额的决定因素}；只满足基本条件但排名不在名额范围内，等于没有。
  \item \textbf{有资格也不等于有学上}——还要过接收院校那一关。
  \item 所以“有资格就够了”是错的，“排名靠前才有用”只说对了一半：排名靠前拿到名额，是必要条件，不是充分条件。
\end{itemize}
\nref{S0454}
\nref{S0639}
\nref{S0568}
\nref{S0205}
\nref{S0131}
\end{qa}

\iin{教育部学生服务与素质发展中心（学信网/研招网）}
\ith{资格认知}
\ist{资格与排名}

\begin{qa}{M0123}{公示前发现成绩录入有误，怎样申请正式更正而不影响推免排名？}
先把“成绩事实”和“推免排名结果”分开处理：向开课学院或教务部门提交成绩复核申请，提供课程、学期、成绩单版本及可核验的原始依据；不要自行改成绩单或要求老师口头改分。取得正式更正记录后，再书面向本院推免工作组申请按更正后的数据重算排名，并要求说明适用的排名截止日期、公式和公示异议期限。推免推荐标准、工作程序由学校制定并公开，异议应走本校申诉渠道；留存提交回执、处理答复和更正前后成绩单。若校内公示期限将到，先在期限内提出书面异议，再补交办理中的证明，避免因等材料错过申诉窗口。
\nref{S0591}
\nref{S0761}
\qscope{操作原则适用于一般情形；具体流程、期限和便利措施以招生单位当年正式通知为准。}
\end{qa}

\section{2.14 规则版本与课程折算}
\begin{qa}{U202610210201}{推免细则在我入学后修订了，低年级学生该按哪一版准备？}
不要把“新发布”直接理解成“对所有年级追溯适用”。推免规则由推荐高校制定并公开，教育部要求招生工作规范、信息公开；但具体实施办法的适用年级、生效时间和过渡安排要看校内文件。先保存本届细则、修订公告和适用对象说明，逐项标出课程范围、成绩算法、资格门槛、名额分配四类变化；再向教务部门确认你所在年级按哪一版执行，并要求指明条款或正式答复。口头答复与网上旧版冲突时，不自行套用对自己有利的口径，等校方书面澄清后按统一规则办理。\nref{S0762}

操作建议： 低年级可以先按当前有效文件建立课程与材料台账，但把尚未确定的年度节点标为“待复核”，在新一届通知发布后更新。
\end{qa}
\begin{qa}{U202610210202}{推免课程从百分制改成等级制，怎样核对成绩折算是否准确？}
先确认推免计算使用的是课程原始等级、学校正式折算分，还是另有推免专用算法；不能自行把“优、良、中”换成百分制后当作官方结果。向教务处申请带有课程名称、学分、原始等级和正式折算口径的成绩记录，同时查本校推免细则是否规定课程清单、计分公式与四舍五入规则。复核时用同一版课程范围和学分权重重算，保留公式、原始成绩单及教务确认；发现课程缺项或折算差异，应在校内公示异议期限内提交书面复核申请。不同学校之间的GPA换算不能替代本校推免排名口径。\nref{S0762}
\end{qa}
\begin{qa}{U202610210203}{推免资格公示后成绩被教务正式更正，排名和名单会自动更新吗？}
不要假定系统会自动联动。成绩更正与推免资格排序通常由不同部门处理：先取得教务部门的正式更正记录，再向学院推免工作组提交更正前后成绩、受影响课程、排名计算表和更正凭证，明确询问是否重算综合成绩、名次及推荐名单。若更正可能改变资格或名次，要求学院说明公示或更正程序、影响范围和处理时间；同时在异议期限内先提交申请并保留回执，不要只依赖电话沟通。最终以推荐高校审核并按规定公示、备案的信息为准。\nref{S0761} \nref{S0762}
\end{qa}
